\documentclass[12pt,a4paper]{article}
\usepackage[T1]{fontenc}
\usepackage[utf8]{inputenc}
\usepackage{textcomp}
\usepackage[ngerman,latin]{babel}
\usepackage{amsmath,amssymb}
\usepackage{geometry}
\geometry{margin=2.5cm}
\usepackage{titlesec}
\usepackage{enumitem}
\usepackage{fancyhdr}
\pagestyle{fancy}
\fancyhf{}
\fancyhead[C]{\small Carl Friedrich Gauß Werke --- Band 5}
\fancyfoot[C]{\thepage}
\renewcommand{\headrulewidth}{0.4pt}

\title{Carl Friedrich Gauß Werke --- Band 5, Teil 2}
\author{Carl Friedrich Gauß}
\date{}

\begin{document}
\maketitle
\thispagestyle{fancy}

\section*{FIGURAE FLUIDORUM IN STATU AEQUILIBRII.}
\addcontentsline{toc}{section}{Figurae fluidorum in statu aequilibrii.}

\setcounter{subsection}{7}
\subsection{}


\[                             -          Bi         --- Idt.cosg.br\]
en
in casu priori
in casu posteriori    = 4r 00 + . dt.cosg.br
rr

denotando
per d? indefinite omnia elementa superficiei spatii s, atque per g, r
eorum respectu eadem, quae antea per literas accentuatas respectu elementorum
determinatorum expressa sunt, denique per x semicircumferentiam eirculi pro
\[radio       ---=   1.\]
Ceterum facile perspieitur, si punctum p. esset neque extra spatium s ne-
que intra, sed in ipsa eius superficie, valere formulam secundam, mutato factore
4r in 2r, siquidem superficies in puncto a neque cuspidem neque aciem offe-
rat; sed ad propositum nostrum haud necessarium est, ad hunc casum attendere.

|                ;                          8.
Per disquisitionem art. praec. evolutio integralis       [fds.dS.p(ds,dS) redu-
citur ad                                                       |
Arnoyo+/fdr.dS. cosy.b(di,dS)
gr asr
si per o denotamus volumen eius spatii, quod utrique spatio s, S commune est,
ita ut prior pars Arob0 excidat, si spatia s, S se invicem excludunt.    Restat
integrale novum, specie etiamnum duplex, revera quintuplex.    Quod ut ad qua-
druplex reducamus, considerabimus integrale
\[                                                  cosg.b(n,dS)\]
fas.       (p, d8)\textregistered\{\}
per omnia elementa spatii S extendendum,                   denotante iterum p punctum deter-
minatum, atque g angulum inter duas rectas ab hoc puncto proficiscentes, alte-
ram versus elementum            dS, alteram fixam.         Hoc integrale, specie simplex,   re-
vera triplex, iam ad aliud integrale revera duplex reducere docebimus, et quidem
duobus modis prorsus diversis.
Planum per punctum p illi rectae fixae normale; per II denotandum, qua-
tenus per proieetionem spatii S attingitur, in elementa infinite parva dII divi-
sum esse concipiatur.    Per punctum talis elementi dIl ducatur recta plano II
normalis, quae deinceps, i. e. progrediendo in directione rectae fixae parallela,
secet superficiem spatii \$ in punctis P', P", P" etc., quorum distantiae a puncto
p. sint resp.         R', R", R''ete.    Similes rectae per omnia puncta peripheriae ele-
V.                                                                           6
42                              PRINCIPIA GENERALIA THEORIAE

menti   dII,    plano ad angulos rectos,       ductae,   spatium prismaticum      formabunt,
et apud puncta P', P", P'' etc. e superficie spatii S elementa exsecabunt, quae
per d7', d7T'', dT'' etc. denotamus. Denique sit y' angulus inter duas rectas
a puncto    P' proficiscentes, alteram extrorsum elemento            dT' normalem, alteram
rectae fixae parallelam, similesque angulos apud puncta              P", P'' etc. exprimant
characteres     y'', y'' etc.    Ita manifesto erit
\[            dll = ---dT..cosy' = +dT''.cosy" = ---dT''.cosy'' = etc.\]

Spatium prismaticum in elementa infinite parva dlII.dz dividatur, denotante z
distantiam puncti indefiniti a plano II (positive acceptam ab ea parte, a qua est
recta fixa); si itaque eiusdem puncti distantiam a puncto x per r designamus,
erit z2= rcosg, nec non (quoniam rr---zz constans est) zdze=rdr,                sive
all.dz.cosqg = dll.dr.    Hine colligitur, integrale nostrum £ds er              )
extensum per eas spatii 8 partes, quae in spatio isto prismatico continentur, ob-
tineri per integrale dII .[Z
ar 2 si extendatur ab r---=R usquad r=R'',
deinab r ---= R'' usque nd r --- R'' etc. Quodsi itaque indefinite statuimus
ano      Re: ii.          :
er


constante integrationis ad lubitum accepta, integrale nostrum pro partibus spatii
S intra spatium prismaticum sitis erit
\[           --- dll.(\$! R---3R'"+9R''---               etc.)\]
\[           --- ---dT'.cosy'.\$R'---dT''.cosy".IR"---dT''.cosy''.YR''---                         etc\]

Collectis his summationibus per prismata omnibus elementis dII respondentia,
manifesto omnia elementa superficiei spatii S exhausta erunt, habebimusque com-
pletum integrale
\[                          fas. =                    --- [aT.cosy.\$R\]
denotante      dT indefinite quodvis elementum superficiei spatii         8, R   eius distan-
tiam a puncto p, atque y angulum inter normalem ad elementum dT' extror-
sum directam atque rectam rectae fixae parallelam.
Hoc itaque modo integrale [fds.dS.p(ds,dS) vodneäue est ad formam
aroyo--- [fat.dT.cosy.d(di,dT)
FIGURAE FLUIDORUM IN STATU AEQUILIBRI.                             43

ubi manifesto      y indicat inclinationem mutuam elementorum                dt, dT,   mensura-
tam per inclinationem normalium utrinque extrorsum respectu spatiorum                         s, S
ductarum,     integrationesque      per superficies completas utriusque spatii extendi
debent.



\subsection*{9.}

Sicuti methodus praecedens divisioni spatii S in elementa prismatica in-
nixa est, ita methodus secunda a divisione eiusdem spatii in elementa pyramida-
lia petetur.   Concipiatur superficies sphaerica radio ---=1 circa centrum g de-
scripta atque in elementa infinite parva divisa. Versus punctum talis elementi
all ducatur a puncto g recta, quae superficiem spatii S secet deinceps in punctis
P', P', P" etc.; distantiae horum punctorum a pn denotentur per R', R'', R"'' etc.
Rectae a p versus omnia puncta peripheriae elementi dII ductae formabunt spa-
tium pyramidale, et apud puncta           P', P", P" etc. e superficie spatii S elementa
exsecabunt,     quae per     d7', d7'', dT''' etc.          designamus.   Denique sit    Q'   angu-
lus inter rectam        P'p atque normalem in elementum dT7' extrorsum ductam, et
perinde sint @'', @'' etc. inclinationes similium normalium apud puncta P'', P'' etc.
ad rectam versus        p ductam.   Ita erit
\[                    >       Treo'    --   --- AT".6Q"                AT'':\]
la                               mar               I zn         ee
valentibus signis superioribus vel inferioribus, prout p. est extra vel intra spa-
tium \$: casus, ubi p est in ipsa superficie spatii S, adnumerandus est casui
priori vel posteriori, prout linea pP' extra vel intra spatium \$ cadit.
Porro patet, pro omnibus partibus spatii \$ intra illud spatium pyramidale
sitis angulum g constantem esse, similique proin modo ut in art. 7 deducimus,
si statuatur indefinite

\[                                     [yr.ar == er\]

constante integrationis ad lubitum accepta, integrale
een                   ds)
\[                                               (p, dS)*\]

extensum per omnes partes spatii S intra illud spatium pyramidale sitas, fore in
casu priori
6 *
44                                   PRINCIPIA GENERALIA THEORIAE

AT'.cosQ'.dR'           dT''.cosQ@''.OR''       dT''.cosQ''AR"''            .
\[            ---=,   608 q . |   R'R'           +       R"R"         ---       R"R"        ==   etc.)\]



in posteriori vero eidem formulae adiiciendum esse terminum dIl.cosg.90.
Iam si haec summatio per omnia superficiei sphaericae elementa colligitur,
integrale completum
\[                                             pen)              ds)\]
\[                                                     (v, 45)\]
fiet
I.     in casu eo, ubi punctum g est extra spatium \$,
\[                                         ---       [    a mie OR\]


denotante         dT   indefinite omnia elementa superficiei spatii S, atque                 Q, R   illo-
rum respectu eadem, quae antea per literas accentuatas respectu elementorum de-
terminatorum expressa sunt, denique g inclinationem rectae a puncto px versus
elementum dT ductae ad rectam nostram fixam.
II. In casu eo, ubi punctum g est intra spatium 8, adiici debet terminus

0.fdll.cosg
ubi q est inclinatio rectae a x versus dII ductae ad rectam fixam, integratioque
per totam superficiem sphaericam extendi debet.     Sed facile perspicietur, inte-
grale istud, extensum per hemisphaerium id, pro quo gq acutus est, fieri = +,
per hemisphaerium alterum autem = ----r; quapropter integrale completum eva-
nescit, valetque pro hoc casu secundo pure eadem formula, quam pro primo tra-
didimus.    Sed aliter. se habet res in casu tertio
III. quoties punctum g. est in superficie ipsa spatii \$8.                  Scilicet hic quoque
adiiciendus est terminus
60. fAll.cosgq
sed integratione per eas tantummodo superficiei sphaericae partes extensa, pro
quibus pars initialis rectae a x versus dII ductae cadit intra spatium \$, sive
(siquidem superficies spatii S in puncto x. neque cuspidem neque aciem offert) pro
quibus haec recta facit angulum obtusum cum recta superficiei spatii \$ in puncto
p normali extrorsumque ducta. Superest itaque, ut integrale hoc sensu acceptum
eruamus.
FIGURAE FLUIDORUM IN STATU AEQUILIBRL.                          45

Secent haec normalis atque recta fixa superficiem sphaericam resp.in punctis
G, H, statuatur arcus GH = k, arcus autem inter @ et punctum indefinitum
superficiei sphaericae =\guillemotright\{\}; denique sit w angulus sphaericus inter arcus \%, v.
Ita erit cosq = cosk.cosv---+ sink.sinv.cosw, et pro dII accipiendum erit ele-
mentum        sinv.dv.dw.        Integrale autem fdll.cosy
\[                     --- /[f(cosk.cose+sink.sinv.cosw)sinv.dv.dw\]

extendi debet       a w=     0    usque ad      w=    360\textdegree\{\},   atque      a v --- 90\textdegree\{\}    usque ad
v ---= 180\textdegree\{\}.        Hoc pacto integratio prior suppeditat
[2rcosk.cosv.sinv.dv

ac dein posterior --- rcosk.
Ad propositum nostrum hic casus tertius eatenus tantum in considerationem
venit, quatenus superficies spatiorum           s, S päartem quandam finitam communem
habent, in qua si punctum g reperitur, eritvel = 0 vel.= 180\textdegree\{\}, adeoque
integrale fall.cosqg vel = ---r vel = +r, prout scilicet apud punctum g
spatia s, S sunt velin eadem plaga vel in plagis oppositis respectu plani utram-
que superficiem tangentis.
Applicando haec ad integrale nostrum primarium ffds.dS.p(ds,dS),                    hu-
ius valor fit
I.     quoties superficies spatiorum      s, S nullam partem finitam communem
habent,
\[                                              d?t.dT.cosg.cosQ.d(dt,dT)\]
\[                          --- Are         + [|        gay\]
II.     quoties superficies spatiorum      s, S partem finitam          ---= 7 communem
habent,
4r00Er700 ge                                    d7)
\[                      =\]


ubi signum superius vel inferius valet, prout spatia s, S sunt ab eadem plaga
vel a plagis oppositis respectu superficiei communis           7.
III.    Quoties superficies spatiorum s, S plures partes finitas discretas com-
munes     habent, sit 7 summa earum, quibus spatia              s, S ab eadem plaga adia-
cent, 7' summa earum, quibus haec spatia a plagis oppositis contigua sunt, erit-
que integrale nostrum
46                           PRINCIPIA   GENERALIA   THEORIAE




Haec tertia formula omnes casus complecti censeri potest.            Integrale duplex per
omnia elementa     utriusque superficiei extendi debet, denotantque          g, @ angulos,
quos facit recta bina elementa di, dT' iungens cum normalibus in haec elementa
extrorsum ductis, directione illius rectae illine a d? versus dT, hinca dT ver-
sus d? accepta.



\subsection*{10.}


Duae transformationes integralis S ds.dS.p(ds,dS)             in artt. 8 et 9 evolu-
tae aequali fere concinnitate se commendant,           proposito autem nostro posterior
magis accommodata est. Problema generale ulterius reduci nequit, nisi ad sup-
positiones determinatas vel circa spatia s, S, vel circa functionem \guillemotright\{\} descenda-
mus.    Et quum functio 9 originem trahat a functione f, disquisitionem ulterio-
rem iam superstruemus eidem hypothesi, a qua ill. LarLace profectus est, puta
vires attractivas moleculares in distantiis insensibilibus tantum sensibiles esse.
Cui phrasi quum aliquid vagi inhaereat, quamdiu non assignatur unitas, ante
omnia observamus,     vim attractivam fr, per functionem distantiae r expressam,
ut cum gravitate    9 homogenea evadat, antea per massam aliquam multiplicari
debere; iam mens illius suppositionis ea est, ut denotante M massam aliquam,
qualis in experimentis occurrere potest, puta quam respectu totius terrae pro ni-
hilo habere licet, Mfr     semper maneat insensibilis respectu gravitatis, quamdiu
r valorem mensuris nostris sensibilem quantumvis parvum habet, dum nihil im-
pediat, quominus valor ipsius Mfr in distantiis insensibilibus non solum sensi-.\_
bilis fieri, sed adeo, decrescente ipsa r,     omnes    limites   superare   possit.   Haud
sane sine admiratione deprehendimus, quam gravia ex hac sola hypothesi, dum
ceteroquin lex functionis fr tamquam omnino incognita spectatur, eruere liceat,
characterem mathematicum prorsus peculiarem prae se ferentia: dum scilicet re-
bus sic stantibus praecisionem mathematicam absolutam sibi vindicare nequeunt,
tamen tantam certissime praecisionem tuentur, ut per nullum experimentum ulla
aberratio a veritate absoluta reperiri possit; quamprimum enim successisset, ta-
lem aberrationem ulli mensurationi subiicere, suppositio ipsa cessaret.


\subsection*{11.}

Supponere licebit, functionem fr (et perinde functionem F'r) attractionem
denotare, omissa ea parte, quae ipsi rr reciproce proportionalis phaenomenis astro-
nomicis explicandis inservit; haec enim pars, quaecungue sit figura fluidi et va-
sis, in quovis puncto insensibilem tantummodo        modificationem gravitati afferre
valet.    Crescente itaque r a valore sensibili in infinitum, /r non modo per se in-
sensibilis erit, sed etiam citius decrescet quam =    Hinc facile colligitur, etiam
integrale f. fr.dr a valore quocunque sensibili in infinitum extensum insensibile
esse, quapropter constantem integrationis fi  fr.dr = ---yr ita acceptam suppo-
nemus, ut habeatur      p00 --- 0, sive ut sit pr ipse valor integralis (fw.dx ab
x =r usquead x = oo extensus.          Hoc pacto pr pro qualibet distantia r de-
notat quantitatem positivam, sed insensibilem, quamdiu r sensibilis est; contra
pro valore insensibili ipsius r non solum sensibilis esse, sed adeo, continuo de-
crescente distantia r, omnes limites superare poterit, sive secundum vulgarem
loquendi modum nihil obstat, quominus sit 90 = ©.


\subsection*{12.}

Inde quod functio gr pro quovis valore sensibili ipsius r insensibilis est et
crescente r continuo decrescit, statim quidem sequitur, integrale frror.dr a
valore aliquo sensibili usque ad alium maiorem extensum etiamnum insensibile
manere, dummodo posterior sit intra ambitum eorum, circa quos experimenta in-
stituere licet: sed neutiquam ex illa proprietate sola concludere fas esset, inte-
grale insensibile manere, ad quantumvis magnum intervallum integratio extenda-
tur. Calculi ill. Lartace ita quidem pronuneiati sunt, ut talem suppositionem
involvant; at dum natura functionis pr incognita est, consultius videtur, ab omni-
bus suppositionibus hypotheticis, quibus supersedere possumus, abstinere.      Quum
itaque constans integrationis [rrpr.dr = ---br arbitrio relicta sit, sufficiat no-
bis, eam ita electam supponere, ut fiat dr = 0 pro valore aliquo sensibili ipsius
r arbitrario, sed intra ambitum dimensionum corporum,          circa quae experimenta
instituere licet.   Hoc pacto dr pro quovis alio eiusmodi valore semper insensibi-
lis erit (positiva pro minori, negativa pro maiori), sed nihil hinc obstat, quomi-
nus pro valore insensibili ipsius r sensibilis evadere possit: addere tamen opor-
tet, phaenomenorum explicationem postulare, ut decrescente distantia r in infini-
tum, valor ipsius dr semper maneat finitus, sive ut d0 sit quantitas finita.      Ce-
48                                   PRINCIPIA GENERALIA THEORIAE
\[              -        cdr                  .                   .                                 .    cthr   ,,-.\]
terum manifesto ---- est quantitas cum gravitate g homogenea, sive \textdegree\{\}?7 Jinea,
adeoque \textdegree\{\}\%\textdegree\{\} Jinea determinata, (pro natura corporum, ad quorum vires attracti-
vas functio fr refertur), cuius magnitudinem ingentem suspicari quidem licet,
sed quam in casibus determinatis vix approximative assignare valemus, saltem non
absque suppositionibus hypotheticis *).                                                    |


\subsection*{13.}

Prorsus simili modo in integratione S dr.dr = ----Hr constantem ita electam
supponimus, ut fiat ör = 0 pro valore arbitrario ipsius r intra ambitum eorum,
pro quibus experimenta instituere licet, quo pacto dr insensibilis erit pro quovis
eiusmodi valore sensibili ipsius r, etiamsi sensibilis evadere possit pro valore in-
sensibili. Manifesto ae exprimit aream figurae duarum dimensionum, adeoque
ie                       .         90      -          ee             ER           .
Er lineam.      Necessario autem 73 est linea magnitudinis insensibilis, quod ita
demonstramus.          Quum       dr inde a r = 0 continuo decrescat, et quidem tam
cito, ut iam insensibilis evaserit, quamprimum r valorem sensibilem acquisivit,
valor ipsius r, pro quo fit dr = \#40, insensibilis esse debet: denotetur ille
\[per p. Consideremus integrale                   [()o---dYr)dr, quodab                r---=0 usquead                r---=R\]
extensum fit = Rb0---90-+BR.    Manifesto hoc integrale maius erit, quam idem
integrale ab  r---=p usquead r= R extensum, atque hoc iterum maius, quam
integrale figo --- bp)dr inter eosdem limites.                         Quare quum integrale postremum
fiat = (0 --- bp) (R---p) = 440.(R---p),                             erit generaliter pro quovis valore ipsius
R (maiori quam p)                                                   Ee
Ry0
\[                                    --- 90 + R>+4\%0.(R---p)\]
Iam si R denotare supponitur valorem fractionis vr, haec relatio suppeditat

\[                                            IR>+Y0.(R---p)\]
quod foret absurdum, si R esset quantitas sensibilis.
\footnote{Concessa   explicatione phaenomenorum      lueis in systemate    emanationis,    refractio   pendet ab at- tractione particularum corporis pellucidi in particulas lucis moleculari,    ratioque refractionis a valore ipsius}
\%o, ita quidem ut habeatur
\[                                                eyo n (nn---1)kk\]
9       sr'
denotante   / longitudinem penduli per minutum secundum vibrantis,          \% motum luminis in vacuo intra mi-
nutum secundum,      n rationem   sinus   anguli incidentiae    ad site anguli refracti:       hoc pacto pro aqua fit
2 bis millies maior quam distantia media solis a terra.
FIGURAE FLUIDORUM IN STATU AEQUILIBRI.                     49

Non obstante itaque ingente magnitudine ipsius        0,   nihil impedit, quo-
minus 00 esse possit quantitas satis modica et cum dimensionibus corporum ex-
perimentis subiectorum comparabilis.


\subsection*{14.}

Superest, ut quae ex hac indole functionis 9 respectu integralis (T)
\[                               er    T.cosg.cosQ.9(dt, dT)\]
\[                                         (dt,d7)*\]

sequuntur, perscrutemur.    Haec investigatio inchoare debet a simpliciori, dum in
alterutra superficie punctum determinatum j consideramus atque integrale (II)
\[                                ee cosg.cosQ.PA(u, di)\]
\[                                        (dt)\]
per totam superficiem ? extendendum evolvimus.           Denotant hie @ angulum in-
ter duas rectas a puncto p proficiscentes, alteram versus elementum dt, alteram
fixam; q vero angulum inter duas rectas a puncto elementi d? proficiscentes, al-
teram versus punetum jr, alteram elemento normalem extrorsumque directam.
Primo loco observamus, si punctum p sit in distantia sensibili a superficie
it, valores omnium d(p, d?) insensibiles fore: in hoc itaque casu totum integrale
(IE) insensibile erit.  Hoc itaque integrale eatenus tantum valorem sensibilem
acequirere potest, quatenus superficies t offert partes in distantia insensibili a
puncto p positas, manifestoque sufficit, integrale (II) per tales partes extendere,
neglectis omnibus, quae sunt in distantiis sensibilibus.
Porro pro in     restituemus ---+dIl, denotante dII in superficie sphae-
rica radio ---=1 circa centrum p descripta elementum id, in quod elementum
d? inde a puncto y visum prolicitur, et valente signo superiori vel inferiori, prout
elementum dt plagam exteriorem vel interiorem puncto p advertit. Hoc pacto
integrale (II) ita exhibetur
\[                               S+all.cosQ.d(p, dt)\]
patetque, huius integralis valorem eatenus tantum sensibilem fieri posse, quate-
nus elementa dII talia, quae ad distantias insensibiles (n, d?) referuntur, spa-
tium magnitudinis sensibilis in superficie sphaerica explent.
Hinc facile colligitur, integrale nostrum, generaliter loquendo, etiam in-
sensibile manere, quoties punctum p. iaceat in ipsa superficie ?: patet enim, pro-
v.                      |                                           7
50                        ..-PRINCIPIA GENERALIA THEORIAE
E                   .



jectiones omnium elementorum        di? a puncto   p insensibiliter remotorum esse in
distantia insensibili a circulo maximo, quem format in superficie sphaerica pla-
'num superficiem t in puncto p tangens.    Excipere oportet tres casus, puta
1) eum, ubi radii curvaturae superficiei * in puncto           px sunt magnitudinis
insensibilis.          |
2) eum, ubi continuitas curvaturae in puncto p, vel intra distantiam in-
sensibilem ab eo interrumpitur (Conf. Disquiss. gen. circa superficies curvas art. 3).
3) eum, ubi superficies ? offert partem aliam a puncto p insensibiliter di-
stantem, puta si apud hoc punctum crassities spatii s est insensibilis.      Ceterum
huncce casum ei, quem in art. seq. tractabimus, adnumerare licet.


\subsection*{15.}

Superest scilicet casus, ubi punctum        non est in superfidie ipsa ?, attamen
in distantia insensibili ab ea: in hoc casu integrale nostrum utique valorem sen-
sibilem habere potest, quem iam accuratius examinabimus.
Secent superficiem sphaericam recta a puncto p normaliter in superficiem
t ducta, atque recta fixa ibinde proficiscens resp. in punctis G, H; statuatur ar-
cus GH=---=\%, arcus autem inter @ atque punctum indefinitum superficiei sphae-
ricae --- v; denique sit w angulus sphaericus inter arcus A, v. Hoc pacto pro
elemento dII accipere licet productum sinv.dv.dw, unde scribendo brevita-
tis caussa r pro (p, d?), integrale (II) fit                     Ä                |
\[                 --- [[+(cosk.cosv+sink.sinv.cosw)Pr.sinv.dv.dw\]

quam integrationem extendere tantummodo oportet per eas partes superficiei
sphaericae, in quas distantiae insensibiles r proiiciuntur. Referuntur
hae ad
partem insensibilem superficiei ?, quam si pro plana habemus, distantiamque
minimam       (puncto   G seu valori v = 0 respondentem) per p denotamus, fit
\[ r=    p\]
z,       Sivea w independens; perfecta itaque integratione respectu variabilis
w, putaa w---         0 usque ad w --- 360\textdegree\{\}, integrale fit
= +fir0r. cosk.cosv.sinv.dv--- pers: potramn
quae integratio extendenda est ab r = o usque ad valorem sensibilem arbitra-
. rium quantumvis paryum. Statuendo itaque generaliter
Auen
.dr \_        ville
r        ua                   Ar.d                        .   i
accepta constante integrationis, ita ut fiat / --- ``--- 0, pro valore arbitrario sen-
sibili intra ambitum eorum, circa quos experimenta instituere licet, erit integrale
(II), neglectis insensibilibus,
= +rcosk.d'p
''


Si dubium videretur, utrum fas sit, partem superficiei ? intra distantiam
insensibilem a puncto p positam pro plana habere, consideremus eius loco sphae-
ricam, et quidem sit R distantia centri sphaerae a puncto p positive vel nega-
tive sumenda, prout centrum est in directione versus G vel in opposita. Ita erit
cp p
\[                               cosv   = Nr\]


\[                         siv.d= Fu---5)---;;]dr\]

unde facile colligitur, integrale pro hoc casu non differre quantitate sensibili a
valore prius invento    ----rcosk.d'p,     simodo     R sit quantitas sensibilis.   Quae-
cunque autem sit curvatura superficiei ? in ea parte, de qua agitur, dummodo
radii curvaturae non sint insensibiles, semper duae superficies sphaericae assignari
poterunt, superficiem. ? in puncto ipsi p proximo tangentes, intra quas £ sita
sit, et quarum radii sint magnitudinis sensibilis, manifestoque tunc integrale
nostrum intra integralia ad illas superficies relata cadet, et proin absque errore
sensibili per eandem formulam exprimetur, quae tunc tantummodo exceptionem
patitur, ubi superficies ? in distantia insensibili a puncto x vel curvaturam radii
insensibilis, vel aciem vel cupidem offert.


\subsection*{16.}

Quodsi iam ab integratione (II) ad integrale (I) progredimur, manifestum
est, hoc insensibile fieri, non solum in eo casu, ubi illa pro nul!o puncto super-
ficiei 7' valorem sensibilem produxit, sed in eo quoque, ubi complexus elemen-
torum superficiei 7', pro quorum punctis integrale (II) sensibile evaserat, aream
tantummodo insensibilis magnitudinis sistit. Quae si rite perpenduntur, appare-
bit, integrale (I) eatenus tantum valorem sensibilem acquirere posse, quatenus
superficies 7' partem vel partes sensibilis magnitudinis contineat in distantia in-
sensibili a superficie ? positas.  Quales partes quum a parallelismo cum superfi-
cie t sensibiliter deviare nequeant, pro quovis earum puncto cosk non sensibi-
7 *
52                                       PRINCIPIA GENERALIA "THEORIAE

liter differet vela +1 vela ---1, prout plaga superficiei 7' exterior vel interior
superficiei 2 advertitur.             Quodsi itaque per t, 7 eas partes superficiei T' deno-
tamus, quae sunt in distantia insensibili a superficie ?, et quidem per \i\{\} eas, ubi
plaga exterior alterius superficiei plagae interiori alterius advertitur, per 7 autem
eas, ubi plagae homonymae sibi mutuo obvertuntur, denique per p distantiam
minimam cuiusvis elementi dr vel dr' a superficie ?, integrale nostrum (I) ne-
glectis insensibilibus fit
\[                                         --- --- [rW'p.dr+ [rh'p.dr\]

Manifesto hic nihil interest, utrum partes t, \i\{\}' ad superficiem                                      T anad t refe-
rantur.
Hoc itaque modo iam nacti sumus solutionem completam problematis, quod
in art. 6 nobis proposueramus, pro ea functionis x indole, cui tamquam basi dis-
quisitio principalis de figura aequilibrii fluidorum innititur, scilicet habemus

\[       [fas.dS.p(ds,dS) = 4nsbo --- 700 +r790 ---rfdr.H' op+rfdr.6'p\]


47.                                72
Origo functionis          9' ita enunciari potest, ut sit
\[                                                H'r ---         [8                     0\]
rr             \guillemotleft\{\}>                                              j

sumto integraliab                  \#---=r usque ad valorem constantem sensibilem arbitrarium,
quem      hic per          R   denotamus.        Manifesto hoc integrale minus erit quam hoc
0)                         FERN:   :
f---
Hr.de       .
,Inter eosdem
or
limites, quod est              2
---= ra       3---,          adeoque a potiori   minus
quam                 Quum autem indefinite habeatur
03.02 dx                   dt2    x                   Vz.dz
\[                               S     8 ee                           |                     F2\}\]
erit



sumto integrali inter eosdem limites, quod minus erit quam integrale f 2                                       ,, ad-
eoque etiam minus quam pr  ---; quocirca valor ipsius ---; maior erit quam
.         .                             \textregistered\{\}              *   ''        8          *        \textregistered\{\}




\[                                                 ---       m    u\]
FIGURAE FEUIDORUM IN STATU AEQUILIBRI.                                 53

Cadit itaque \#r inter limites
BR     ?
dr atque      Hr.                      r

quorum differentia, decrescente r in infinitum, manifesto quavis quantitate as-
signabili minor evadere potest, quum supponamus esse vel d0 quantitatem fini-
tam.     Colligimus hinc, statui debere           0'0 --- \#0.        Patet itaque, in formula, ad
quam      in art. praec.     pervenimus,       terminum           ---r700      tamquam     sub termino
--- rfdr.0p,       atque terminum         7700     tamquam sub termino             x far.dp    compre-
hensum considerari posse, si distinctionem,                 quam    inter distantiam     insensibilem
et distantiam nullam fecimus, tolleremus, atque partes 7, 7' resp. partibus r, 7'
adnumeraremus. Sed gquamquam hoc modo solutionis elegantia sensu mathematico
augeretur, tamen ad propositum nostrum praestat, distinctionem illam conservare.


\subsection*{18.}

In applicatione disquisitionis praecedentis ad evolutionem termini secundi
expressionis 2 art.3,        spatium
secundum inde ab art. 6 per S denotatum cum
primo identicum est; quae itaque in art.16 erant 0, 7,7', hic erunt s, t,0, si
ti denotat totam superficiem spatii s a fluido impleti.                     Quapropter   quoties hoc
spatium neque partes sensibilis extensionis sed insensibilis crassitiei continet, ne-
que eiusmodi interstitia (fissuras), pars secunda expressionis © fit
.                                   --- +rcc(sb0 --- 100)

,Exceptiones itaque adsunt duae:
1) Si spatium s continet partem insensibilis crassitiei, huius superficies
duas partes sensibiliter aequales offeret, quarum alterutra per ?' denotata, cras-
sitieque spatii apud quodvis elementum dt' indefinite per p, accedet expressioni
praecedenti terminus

rec/hp.dt'
2) Si spatium      s continet cavitatem insensibilis crassitiei, accedet similis
terminus, puta zce/\textregistered\{\}p.dt'',             denotante         ?'' alterutram partem superficiei t fis-
surae contiguam, atque p indefinite crassitiem fissurae in quovis puncto.
In evolutiong termini tertii expressionis @ signum S retinendum erit, ut
denotet spatium a vase repletum, sed loco characteristicae f characteristicam F
54                           PRINCIPIA    GENERALIA    THEORIAE


ad vim attractivam molecularum vasis relatam substituere oportebit, et perinde
loco functionum per characteristicas 9, 4, 0, \#" denotatarum alias per characte-
risticas \textregistered\{\}, W, 0, 0\textdegree\{\} denotandas adhibere, quas perinde ab F' pendere supponi-
mus ut illas ab f. Quae in disquisitione generali erant 0,7', hic manifesto erunt
0: pro 7 vero hic simpliciter literam 7' adoptabimus, ut indicet non superficiem
totam spatii S, sed eam partem, quae fluido contigua est. Hoc pacto pars ter-
tia expressionis   © fit, generaliter loquendo,

\[                                      ---    zcCT90\]

exceptis etiam hic duobus casibus, puta
3) Siapud partem sensibilem 7'' superficiei                7' fluidum crassitiem insensi-
bilem habet, indefinite per p exprimendam, accedet terminus
\[                                     ---rcCf[\$p.dT'\]
4) Si superficies vasis praeter partem         T fluido contiguam, offert aliam 7''
in distantia quidem sed insensibili a fluido positam, accedet, denotante p inde-
finite hanc distantiam pro quolibet puncto, terminus
\[                                     +rcC[p.dT''\]
Superfluum foret, exceptioni primae, quatenus sub tertia non continetur, nec non
secundae vel quartae immorari:       etiamsi enim aequilibrium fluidi in casibus qui-
busdam huc referendis,     attamen    maxime     specialibus,        locum habere queat, tale
aequilibrium nec stabile neque experimentis accessibile esse posset.              Contra casus
exceptus primus, quatenus sub tertio continetur, utique theoriae essentialis .est,
verumtamen aliquantisper hic seponetur, ut conditiones aequilibrii, quatenus abs-
que cute fluidi insensibili, vasi adhaerente, consistere potest, explorentur.
Dum itaque omnes has exceptiones seponimus, expressio, cuius valor in
statu aequilibrii maximum esse debet, haec erit

---9 cfzds+\$cesyo           --- rec              +rceC T90

et quum   in omnibus    mutationibus,      quas figura fluidi subire potest, spatium s
invariatum maneat,     expressio sequens

Sz4s +7.                     Ele

in statu aequilibrii minimum esse debebit.
FIGURAE FLUIDORUM IN STATÜ AEQUILIBRIL.                           55

Jam supra monuimus, vu exhibere spatium duarum dimensionum,                    idem-
que de -         valet.   Statuendo itaque



erunt       a, 6 lineae constantes a relatione gravitatis ad intensitatem virium,           quas
,partes fluidi a se mutuo et a moleculis vasis patiuntur, pendentes; et si porro
partem liberam superficiei fluidi, i.e. eam, quae vasi non est contigua, per U de-
notamus, ut habeatur \{= TU,       minimum esse debet in statu ,aequihbrüi eX-
pressio sequens, abhinc per W denotanda:

Szds+ (aa --- 26           TtaeU



\subsection*{19.}

Antequam quae ex hoc theoremate gravissimo sequuntur generaliter et com-
plete evolvamus,  operae pretium erit ostendere,               quanta facilitate phaenomenon
principale tuborum capillarium inde demanet.
Consideremus fluidum in aequilibrio in vase bicrurali, ita ut pars superfi-
ciei liberae fluidi sit in primo crure,        pars alia in secundo:     parietes vasis in con-.
finiis-harum partium verticales supponimus.                Sit a area sectionis horizontalis in-
ternae primi cruris (vel exactius proiectionis horizontalis superficiei liberae fluidi
in primo crure), 5 eiusdem peripheria, denique ah volumen fluidi in hoc crure,
pariete verticali deorsum usque ad planum, a quo numerantur distantiae z, con-
tinuato, sive, quod eodem redit, A altitudo media fluidi supra hoc planum: simi-
lia denotentur pro secundo crure per literas a', b, 4.              Si statum fluidi mutatio-
nem infinite parvam subire concipimus,            et quidem talem, ut utraque superficiei
liberae pars figuram       suam    servet,   variatio partis primae expressionis      W,   puta
integralis fizds,      manifesto erit

\[                                        ---= ahdh+ ahdh\]
variatio ipsius 7 autem                                |
\[        =                               --- bdh+bdh\]
denique per hyp. dU = 0. Hinc colligitur

\[                    dW = ahdh---+ahdh---(266--- ac) (bdh---+bdah)\]
56                                    PRINCIPIA GENERALIA THEORIAE

Porro quum volumen integrum fluidi invariatum maneat, erit
\[                                            adh-+adh =        0\]
et proin                                                                         :
\[                        aW = dhla(h---\]
\[                               N) ---(286--- aa) --- ``\%)]                        [77\]


Conditio itaque, ut W in statu aequilibrii sit minimum, perducit ad aequatio-
nem, phaenomenon principale tuborum capillarium implicantem
\[                                     h---k = (266---aa)(---5)       [7\]



sponteque patet, huic aequationi revera respondere valorem minimum ipsius W,
.  .  adWw              aa  .                 ea
quum valor ipsius dh Hat --- aA , 1. e. natura sua positivus.
a'      .
Crus secundum priori largius pronunciatur, si quotiens 7 est malor quam
re fluidum itaque in crure arctiori magis depressum vel magis elevatum erit quam
in largiori, prout quadratum 65 minus vel maius est quam 4aa; et si forte ha-
beretur    65 = 4aa,              altitudo in utroque crure eadem foret.             Si crus secundum
d'         3              b
tam largum est, ut           =     negligi possit prae   7;   erit prox\i\{\}me
b
\[                                           h---k=     2856 --- aa)---\]
a


In tubis itaque capillaribus cylindricis fluidi depressio vel elevatio diametro tubo
reciproce proportionalis est.              Haec omnia tum cum experientia tum cum iis, quae
ill. Larptace per theoriam stabilire conatus est, conveniunt.
Si vas pluribus cruribus verticalibus inter se communicantibus instructum
est, designent       a', b', A'       pro tertio,   a'', b'', A''     pro quarto etc. eadem,        quae
a,b, h pro primo, eritque etiam

\[                                    h---h" = (266---aa) (2----,)\]
\[                                    h---h" = (266---aa)(?---,)\]
Conecinnius hae aequationes ita exhibentur:
b       r                     f                       n         a               m
\[h--- (2656---aa)_ --- h--- (266 ---aa)-, = W--- (28586--- u), --- h" --- (266 ---aa)>,ete.\]

Quum planum horizontale, a quo altitudines numerantur,                       arbitrarium sit, patet,
si illud ita assumatur,           ut sit
\[                                            h= (255---oo)\]
FIGURAE FLUIDORUM IN STATU AEQUILIBRU.                              57

etiam in reliquis cruribus fore

\[     "= (266---oao)-,          A'= (266---.aa) aa\]
m

\[                                                                    (265 --- ao)a'' etc.\]

Hocce planum, cuius conceptum infra generalius stabiliemus, vocari potest pla-
num horizontale normale (plan de niveau). Supponendo (si opus est) parietes ver-
ticales singulorum crurium usque       ad hoc planum        productos,    ah, ah, a'h" etc.
expriment, pro 2665>oa, quantitates fluidi in singulis cruribus supra hoc pla-
num elevati, vel pro 2665<aa, quantitates fluidi infra hoc planum in singulis
cruribus deficientis: hae itaque quantitates aequales sunt productis ex area con-
stante   26565---oa   in circumferentias     b, D', 5", b'' etc.


\subsection*{20.}

Superest iam, ut e theoremate art. 18 indolem figurae aequilibrii determine-
mus, cuius negotii cardo vertitur in evolutione generali variationis, quam ex-
pressio W patitur, dum figura spatii a fluido impleti mutationem quamcunque
infinite parvam subit. Sed quum calculus variationum integralium duplicium pro
casu, ubi etiam limites tamquam variabiles spectari debent, hactenus parum ex-
cultus sit, hanc disquisitionem subtilem paullo profundius petere oportet.
Considerabimus superficiei, quae spatium s a reliquo spatio separat, par-
tem U, atque quodvis illius punctum per tres coordinatas x, y, z determinari
supponemus, quarum tertia sit distantia a plano horizontali arbitrario.    Spectari
itaque poterit z tamquam functio indeterminatarum \#, y, cuius differentialia
partialia secundum morem suetum, sed omissis vinculis, per
dz             dz
j,.de,         a, dy



denotabimus.    In quovis superficiei puncto rectam superficiei normalem et re-
spectu spatii s extrorsum directam concipimus, cosinusque angulorum inter hanc
normalem atque rectas axibus coordinatarum             @, y,z    parallelas per   \&,n,C   de-
notamus.    Hoc pacto erit

\[                                  E+nn+\%              = 1\]
dz AN. E3           dz
RE          |
Y Braga S'          Zu         L

Limes superficiei U erit linea in se rediens, quam per P denotamus, et dum
v.                                                         8
58                          PRINCIPIA   GENERALIA    THEORIAE


motu continuo descripta supponitur, eius elementa dP (perinde ut elementa su-
perficiei dU) semper positive accipiemus.     Cosinus angulorum, quos directio
elementi dP facit cum axibus coordinatarum w, y,z, per X, Y, Z denotamus:
ne vero sensus directionis ambiguus maneat, hanc ita decernimus, ut ipsa primo
loco, directio normalis in elementum dP superficiem U' tangentis et huius re-
spectu introrsum ductae secundo loco, denique normalis in superficiem respectu-
que spatii s extrorsum ducta tertio loco, constituant systema trium rectarum si-
militer deinceps sitarum, ut axes coordinatarum w, y, z. Ita facile perspicietur
(cf. Disquiss. gen. circa superficies curvas art.2), cosinus angulorum inter directio-
nem illam secundam atque axes coordinatarum \textregistered\{\}, y, 2 esse resp.
\[                    W"Z---CY,       OX---EZ,          EY---\i\{\}rX\]
si \&\textdegree\{\}, m, C\textdegree\{\} sint valores ipsarum \&, n, \& pro puncto elementi         dP.

21...
His ita praeparatis supponamus,       superficiem     U pati mutationem   qualem-
cunque infinite parvam.    Si sufficeret, tales tantummodöd mutationes considerare,
pro quibus limes P semper invariatus, vel saltem in eadem superficie verticali
maneret, manifesto soli coordinatae tertiae z variationem inducere oporteret, quo
pacto problema longe facilius evaderet; sed quum problema maxima generalitate
nobis ventilandum sit, in tali investigationis modo consideratio variabilitatis limi-
tum in ambages incommodas concinnitatemque turbantes perduceret; quamobrem
praestabit, statim ab initio omnes tres coordinatas variationi subiicere.       Rem
itaque sic imaginabimur, ut cuivis puncto superficiei, cuius coordinatae sunt
x, y, 2, substituamus aliud, cuius coordinatae sint \&+0@, y+öy, z2+6z,              ubi
6\%, 6y, 62 spectari possunt tamquam functiones indeterminatae ipsarum z, y,
sed quarum valores manent infinite parvae. Inquiramus nunc in variationes sin-
gulorum elementorum expressionis         W, et quidem initium faciamus a variatione
ipsius elementi dU.
Concipiamus elementum superficiei U triangulare dU              inter puncta, quo-
rum coordinatae sint

Ve,         Y;             P4




ade,        y+dy,          de up

Area duplex huius trianguli per principia nota invenitur
\[                    er (da.dy---dy.da)yf\i\{\}+()                   +)dy ]\]
si, quod licet, supponimus, de.d'y---dy.d'x esse quantitatem positivam.
In superficie variata loco illorum punctorum tria alia habebimus, quorum
coordinatae erunt
puncti primi             z-+ö@,         y+öoy,      z+6z

puncti secundi

x+de+öc+T\%.de +: dy
yHay+öy+T: Ant BU,
+...da+, .dy-+öz AH de +                        dy
puncti .-tertii
ddx        ddz     y
\[                  a +de +64... deer                  .dy\]

yhayhöy rin.dar .dy
dz   n     =2 dyhösh.
de ddz
\[                                                               --- d'y\]
pr

ER duplex trianguli inter haec puncta invenitur per eandem methodum
\[                               = (dae.dy---dy.da)/N\]
si brevitatis caussa per N denotatur aggregatum

Un)        dy
GESSUK FÜ TAG Zar Aa
et          latae             mn ha)]
Facta evolutione et reiectis quantitatibus secundi ordinis, invenitur

N
Vet]                                   nn
sr Sl
si brevitatis gratia per L denotatur aggregatum
\[                                                                       g*+\]
60                                 PRINCIPIA      GENERALIA         THEORIAE

ddr        dz,2   döxz   dz   dz       döy   dz    dz          ddy                          ddz          dödz    dz
ot                oo                   SEeRTESN yaDH:                                             tt2,
Est itaque ratio trianguli primi ad secundum ut 1 ad
We L
ar,        de           dy

adeoque independens a figura trianguli                  dU,         resultatque

\[                                     ddU---                     hs\]
de.''         dg   .*
\[                                        |          1+(7,)             a2)\]
sive in terminis explicitis -

SAU
\[  = AUT                                 En                     + EHE                                                 end)\]
dy



\subsection*{22.}

Variationem totius superficiei U obtinebimus per integrationem huius ex-
pressionis per omnia elementa dU extendendam.      Ad hunc finem duas huius in-
tegralis partes, puta
\[                       fav.( mn                                       --- KT)             =    4\]

atque
SAU,
\[                         +Hei+                                                      en) =\]
seorsim tractabimus.
Concipiatur planum axi coordinatarum                     y normale,              et quidem tale, ut va-
lor determinatus ipsius y, ei competens, sit intra ambitum valorum extremorum,
quos habet y in superficie U. Hoc planum peripheriam P secabit velin duo-
bus, vel in auener vel in sex etc. punctis, quorum coordinatae primae sint de-
inceps ='', \guillemotleft\{\}', \&'' etc.; perinde reliquae quantitates ad haec puncta pertinentes per
indices ER                    Eodem modo secetur superficies per aliud planum illiin-
finite propinquum et parallelum,             cui competat coordinata secunda y-+-dy;                                  inter
haec plana reperientur elementa peripheriae                         dP\textregistered\{\}, dP'', dP'' etc.,               perspicietur-
que facile, haberi

\[             dy=---Y'daP'=+Y'daP'                        = ---Y'aP"                     = Y''dP''            etc.\]

Si insuper concipimus infinite multa plana axi coordinatarum            \# normalia,    cuivis
elemento d\# inter ='' et \guillemotleft\{\}', velinter         \guillemotleft\{\}'' et \textregistered\{\}'' etc. sito respondebit elementum
au = u       , unde patet, eam partem integralis A, quae respondet parti super-
ficiei inter plana y, y+dy    sitae, haberi ex integratione
Bert ddr in döy\_              ,,dd
ayfaw.          ABT   .ErTE              a;)
extensaab    2 = 2\textdegree\{\} usquead        2 ---=.\guillemotleft\{\},    denab       = x'' usquead           \# = 7'' etc.
Indefinite vero hoc integrale exhibetur per
eo                  en
a'?        .
ER            N. \&y--- töndy---dyfien,                   ----    öy. ---        d2.2)dx

unde colligitur, prodire pro casu nostro
iR    SD   Ö \_L.öy\textdegree\{\}---£*8\textdegree\{\})Y!aP\textregistered\{\}

--- tr            .6W a        .öy ---E0z7)Y'dP'
er               a         09       Erb Frar?
\[                  --- etc.\]
en          R
\[                  --- [CAU. (dw.         \_\_ ---dy..2---62.5;)\]

sive, quod idem est,
anti               in
zer        gr ---N.8y--- E82) FAP--- fLaD. (ir.                    ------ ---dy.
\[                                                                    2 ---\%.E)\]
ubi tum summatio per omnia elementa          dP,   tum integratio per omnia elementa
dU, intra plana y et y+dy         sita, extendenda est.

Tota itaque quantitas A exprimetur per
ann    8\%          a
SEHE. 30 ---2.8y---E82) YAP---[CAU.(ö
ubi integrationem priorem per totam peripheriam          P, posteriorem per totam su
perficiem U extendere oportet.
62                                   PRINCIPIA   GENERALIA           THEORIAE



\subsection*{23.}

Per ratiocinia prorsus similia invenimus
en       BE+EL        dn
\&          SEH+LL                                                    L         c
B =      [(7.5\textdegree\{\}---       7       .öy+nd2)XdaP+[ldU(dr. >, ---6y.                             5      +d2.4,)

Statuendo itaque, pro quovis puncto peripheriae P,

\[                                              FEl)dez = (IQ\]
XEn + Yan +L9))de - [X EE+LC) + Yenlöy+Xnt---
nec non, pro quovis puncto superficiei                 U,


\[             dar                                                           tler) =\]
erit tandem

\[                                       SU=/QdP+/[vau\]
ubi integratio prima per totam peripheriam P, secunda per totam superficiem                                U
extendi debet.


\subsection*{24.}

Formulas pro Q@ et V modo allatas notabiliter contrahere licet. Et quidem,
adiumento aequationis X\&+-Yn+ZT = 0, Q statim induit formam symme-
tricam sequentem:

\[                   Q = (YC--- Zn)de + (ZE---XL)6y+(XKn---Y\})6z\]

Quo etiam expressio pro V eruta in formam concinniorem reducatur, ob-
servamus,      e formulis

ds \_ 0                    Se
de          BE                 ang    c
sequi
\&              N
ei
dy    rn,

Hinc fit
in                          E                        n
\[                             ---._                       "r                  Bet\]
FIGURAE    FLUIDORUM      IN STATU     AEQUILIBRI,                     63
\&
Porro ex     E\&\&+nn+\{\{ =         1 deducimus
de        BET              GER

atque hinc
Fikuid             47                               aa
ERENTOER
SE          WMILT          CAT                      I x
Substitutis his valoribus in coöfficiente ipsius 6x in expressione pro                   V, ille fit
2
\[                                           -te+m\]
d\& , dn

Prorsus simili modo coöfficiens ipsius öy in eadem expressione transit in
d\& | dn
it
Hoc itaque pacto nanciscimur
|                                                     .
\[                             =       Fr      tny           ++)\]


f                             25.
Antequam ulterius progrediamur, significationem geometricam expressio-
num erutarum illustrare conveniet.      Ad hunc finem directiones varias hic occur-
rentes intuitioni faciliori subiiciemus sequendo eum modum, quem in Disquiss.
gen. circa superficies curvas introduximus, puta referendo illas ad puncta super-
ficiei sphaericae radio ---= 1 circa centrum arbitrarium descriptae.    Primo itaque
directiones axium coordinatarum \#, y, z denotabimus per puncta (1), (2), (3);
dein directionem normalis in superficiem et respectu spatii s extrorsum ductae per
punctum (4); demique directionem rectae a quolibet superficiei puncto versus
ipsius loecum variatum ductae, per punctum (5). Variationem loci ipsam, seu
quantitatem Y(62\textdegree\{\}-+0y?-+62\textdegree\{\}?),               semper positive sumendam, brevitatis caussa
per Öe denotabimus, arcumque inter duo sphaerae puncta, ute.g. (1) et (5),
sive angulum, qui illum arcum mensurat, ita (1, 5) scribemus. Erit itaque
ör ---= öde.cos(1,5),            öy= de.cos(2,5),,                Öz--- Öe.cos(3, 5)
Haec pro quovis superficiei puncto valent. In eius limite, seu peripheria P,
duae aliae directiones accedunt.   Primo directio elementi dP, cui respondeat
64                                PRINCIPIA GENERALIA THEORIAE

punctum (6); dein directio rectae huic normalis superficiem tangentis eiusque re-
spectu introrsum ductae, cui respondeat punctum (7). Per hypothesin nostram
puncta (6), (7), (4) eodem ordine iacent, ut (1), (2), (3), observetur praeterea,
(4,6), (4,7), (6,7) exhibere quadrantes seu angulos rectos.     Ita prodeunt ae-
quationes iam supra (art. 20) traditae

\[     n2---CY= cos(1,7),                   I[X---EZ\]
\[                                           = c0s(2,7),                           EY---yX\]
\[                                                                                    = cos(3,7)\]

formulaeque art. praec. has formas induunt:

Q = ---öe.cos(5,7)
\[                            V=     de.cos(4,5). Et)\]

Exprimit itaque @ translationem cuiusvis puncti peripheriae P a plano hanc tan-
gente superficiei U normali, in plaga ab hac aversa positive sumendam; factor
ipsius V autem Öe.cos(4,5) manifesto indicat translationem cuiusvis puncti su-
perficiei U a plano hanc tangente, positive sumendam in plaga a spatio s aversa.
Sed etiam factorem alterum ipsius V per significationem geometricam ex-
plicare licet. Habemus enim

\[                             i=---i..dz                               ar\]
0.dz
1                  dz,?               dz,?
air                           trt)
Hinc prodit


m m
\[                                   \        je? [=?]LS\}\]
EN
fe? u
I


ddz   ddz
\[                                         ---c(ny             de ,+Rz. dy\]
A             Ads                        1:         ddz
dy          Br,                           ten      rer dy

---E ieepege Er
en ddz
et proin
de,   an.             ddz           dz 2               2ddz    d       dd   d
\[            ut,         =           1+i4,)                I--- de.\]
TEE 12\textdegree\{\} rl                 ital
FIGURAE FLUIDORUM IN STATU AEQUILIBRIL.                      65

cuius expressionis valorem constat esse
i      1
\[                                      =,;+t75\]
denotantibus   R, R' radios curvaturae extremos in puncto de quo agitur, et qui-
dem positive accipiendos,   quoties convexitas superficiei extrorsum vertitur.



\subsection*{26.}

Examen attentum analysis nostrae inde ab art. 22 patefaciet suppositionem
tacitam illi adhaerentem, scilicet quibusvis valoribus coordinatarum x, y unicum
tantummodo valorem ipsius 2 respondere, atque valorem ipsius \{ ubique per
totam superficiem U esse positivum. Nihilominus veritas theorematis finalis, ad
quod analysis ista perduxit, puta (I)                |
\$U---= ---fde.cos(5, 7).dP+föe.cos(4,5).(5+75,)AU

ad hanc suppositionem non restringitur, sed generaliter valet. Quam generalita-
tem si statim ab initio amplecti voluissemus, vel quasdam ambages incurrere, vel
methodum aliquantum diversam sequi oportuisset: sed ad eundem finem etiam
per considerationes sequentes facile pervenire licet.
Analysis nostra manifesto independens est a suppositione, quod axis coor-
dinatarum z est verticalis, quin potius in illa situs axium prorsus arbitrarius ma-
net, veritasque theorematis stabilita est pro omnibus superficiebus, pro quibus
complexus omnium punctorum (4) unico hemisphaerio includi potest; sufficit enim,
talis hemisphaerii centrum (polum) pro (3) adoptare.
Si vero proponitur superficies huic conditioni non satisfaciens, certe in duas
pluresve partes dispesci poterit, quae singulae tali conditioni satisfaciant.    Iam
facile perspicietur, si superficies quaedam in duas partes divisa fuerit, veritatem
theorematis pro figura tota statim sequi e veritate pro singulis partibus.   Constet
enim figura U e partibus U', U'', sitque P' peripheria figurae U', atque P"
peripheria figurae U''; porro habeant P', P" partem communem P'', ita ut P'
constet ex P'' et P'', P" veroex       P'' et P"'', unde manifesto peripheria figurae
U integra P constabit ex P"'' et P"''.          Ita erit quidem
P"''
fe .
cos (5, 7)dP' = föe. cos(5,7)d P"+föe.cos(5,7)d
\[                                fe. cos ( les                   cos (gps\]
\[          fe .cos (5, 7) ap.\]
v.                                                                   9
66                                PRINCIPIA GENERALIA THEORIAE

sed probe notandum, valorem integralis f de.cos(5,7)AP'', quatenus est pars
prioris integralis, exacte oppositum esse valori eiusdem integralis, quatenus est
pars posterioris integralis, quum cuivis puncto lineae P'', in his duobus casibus
directionibus oppositis describendae,          loca puncti (7) opposita adeoque valores op-
positi factoris     cos(5,7)    respondeant.     In additione itaque hae partes sese de-
struunt, fitque
föe.cos (5, 7)AP''+föe.cos(5,7)d.P" --- [öe.cos(5,7)dP

unde, quum habeatur \&U --- 6U'+6U", valor ipsius 6U cum formula allata
(I) conspirans sponte demanat, dum haec formula cum valoribus variationum
8U', 6U'' quadrare supponitur.
Denique observamus, veritatem theorematis (I) etiam e considerationibus
geometricis hauriri potuisse, et quidem facilius quam per methodum analyticam,
quam tamen hic ideo praetulimus, ut occasionem, calculo variationum, pro inte-
gralibus duplieibus limitibus variabilibus inclusis parum hactenus exculto,, ali-
quid lucis effundendi arriperemus, methodum alteram geometricam satis obviam
lectori perito relinquentes.


\subsection*{27.}

Superest, ut variationes evolvamus, quas elementa reliqua expressionis W
per variationem figurae spatii s patiuntur, et primo de variatione voluminis spa-
tii s agemus.
Resumamus        duo triangula in art. 21 considerata,          iungamusque laterum
puncta respondentia per rectas, ut oriatur solidum, cuius loco accipere licet prisma
basis AU, altitudinis          Edo-ndy-+Löz = Öe.cos(4,5),            et quidem haec forma
dabit altitudinem in forma positiva seu negativa, prout triangulum transpositum
et proin totum solidum iacet extra vel intra spatium s.              Hinc habemus (II)
\[                                     os --- [AU.de.cos\]
\[                                                   (4, 5)\]

Porro hinc sequitur, variationem integralis [zds esse (III)
öfzds --- [zdU.de.cos(4, 5)

Quod       vero attinet ad variationem     quantitatis    7', ante omnia observamus,
quum P denotet limitem communem superficierum                   7, U, transpositiones puncto-

rum peripheriae P satisfacere debere huic conditioni, ut loca nova in superficie
spatii S maneant. Manifesto itaque per transpositionem elementi dP, super-
ficies T patitur mutationem +dP.6e.sin(5, 6), perspicieturque facile, genera-
liter loquendo signum positivum vel negativum a signo quantitatis cos(4,5) pen-
dere.    Sed concinnius haec variatio exprimitur introducendo directionem no-
vam, quae sit in plano superficiem spatii S tangente, lineae P normalis, et re-
''

spectu spatii s extrorsum ducta.    Denotando per (8) punctum huic directioni
respondens, variatio superficiei    7' a transpositione elementi           dP     oriunda erit
dP.Öe.cos(5, 8), sive (IV)                    |
\textdegree\{\}T --- [dP.öe.cos(5,8)

ubi signum factoris   cos(5, 8) sponte decidet,    utrum mutatio sit incrementum an
decrementum.
Quum punctum         (6) sit polus circuli maximi            per puncta        (7), (8)
(5) iaceat in circulo maximo
ducti, punetumque                                          per puncta       (6), (8)       ducto,
,puncta (5), (7), (8) formabunt triangulum in (8) rectangulum, eritque adeo
cos(5,7) = cos (5, 8).cos(7,8) : arcus (7, 8) autem est mensura anguli inter duo
plana superficies spatiorum s, S in eorum intersectione P tangentia, et quidem
inter eas horum planorum plagas, quae spatium vacuum               includunt.       Hunc an-
gulum per ö denotabimus, unde 180'---i erit angulus inter planorum plagas
eas, quae spatium s continent, formulaque nostra (V)
\[                               cos (5, 7) = cos (5, 8). cos.\]


\subsection*{28.}

E combinatione formularum      I....IV      prodit variatio expressionis         W

\[                     5W = [AU.de.cos (4,5). g+aa(5+z7)]\]
--- [a P.de.cos(5, 8). (aacosi---
aa 266)

ubi integrale prius extendi debet per omnia elementa         dU     partis liberae super-
(si forte plures separatae adsint), integrale
ficiei spatii s, vel partium liberarum
posterius autem per omnia elementa d.P lineae vel linearum, quae illam partem
liberam, vel illas partes liberas a reliquis spatio S contiguis separant.
Iam quum in statu aequilibrii valor ipsius W debeat esse minimum,
"adeogue   admittere   nequeat   mutationem     negativam     pro    ulla    mutatione       infi-
9, *
68                          PRINCIPIA GENERALIA THEORIAE

nite parva figurae fluidi, pro qua volumen       s invariatum manet, i. e, pro qua
Os F dU.öe.cos(4,5) evanescit, facile perspicietur, figuram superficiei U in
statu aequilibrii talem esse debere, ut in omnibus eius punctis elementum varia-
tionis 6W     hoc
\[                          dU.oe.cos(4,5). @-+aa(5-+ 59)]\]

proportionale sit elemento variationis 6s, puta quantitati dU.Ööe.cos(4,5),     sive
quod idem est, ut fiat

\[                             z+aa(54+7) = Const.\]

Manifesto enim, si haec proportionalitas locum non haberet, valor ipsius W de-
crementi capax foret per idoneam mutationem figurae superficiei U, limite P
adeo invariato manente.    ÜCeterum aequatio illa pro Zota superficie U valet,
etiamsi haec e pluribus partibus separatis constet,     dummodo   fluidum ipsum co-
haereat.
Aequatio ista constituit theorema fundamentale primum in theoria aequili-
brii luidorum, quod iam ab ill. LarLace erutum est, sed per methodum a nostra
plane diversam.
Si planum, pro quo 2 quantitati aequationis constanti aequalis est, et quod
planum horizontale normale (plan de niveau) vocare possumus,       loco eius, a quo
coordinatae   2 numeratae erant,   adoptamus,    erit
1      1
\[                                   =   Res)\]


unde protinus demanant corollaria sequentia.
I. Si planum normale superficiem liberam U ullibi secat, in quovis sectio-
nis puncto superficies necessario concavo-convexa erit, atque radius maximus
convexitatis radio maximo concavitatis aequalis.
Il. Supra planum normale superficies vel concavo-concava erit, vel, sicubi
fuerit concavo-convexa,    curvatura concava convexam superabit.
III.   Infra planum normale superficies vel erit convexo-convexa, vel, sicubi
fuerit concavo-convexa, curvatura convexa concavam superabit.
IV. Superficies libera U nequit habere partem finitam planam nisi hori-
zontalem et cum plano normali coincidentem.
FIGURAE FLUIDORUM IN STATU AEQUILIBRI.                            69


\subsection*{29.}

Aequatione, quam modo stabilivimus, subsistente, variatio valoris ipsius W
reducitur ad                                                              |

2566)
,W = ---[dP.öde.cos(5,8) (aa cosi---aa+
\_   unde introducendo angulum A talem ut sit
\[                        cos    A ee   ---.\]
\[                                         oa     sive        in44=-\]
habemus
\[                                        8). (cos A--- cos)\]
\textdegree\{\}W = aafdP.öe.cos(5,

integratione per totam lineam P extensa.     Memores esse debemus, factorem
cos(5, 8) aequalem esse ipsi sin(5,6), signo positivo vel negativo affecto, prout
fluidum in motu suo virtuali apud elementum                dP    vel ultra limitem   P redun-
dare, vel citra recedere concipitur.      Hinc facile concludimus, in statu aequilibrii,
generaliter loquendo, ubique esse debere i---= A. Si eenim in aliqua parte lineae
P esset i<(A, motus virtualis primi generis in hac parte, manente parte reli-
qua limitis P invariata,      manifesto ipsi        W   variationem negativam induceret,      et
perinde negativa variatio ipsius      W   prodiret per motum virtualem fluidi secundi
generis, siin ulla parte lineae P esset i>A: utraque igitur suppositio condi-
tioni minimi in aequilibrio adversatur.
Hoc est theorema       fundamentale      secundum,        quod etiam    investigationibus
ill. Laptace intertextum, sed e principio virium molecularium haud demonstra-
tum videmus.



\subsection*{30.}

Theorema art. praec. modificatione quadam eget in casu singulari, quem si-
lentio praeterire non licet. Tacite scilicet supposuimus, superficiem vasis iuxta
totum limitem P curvatura continua gaudere, ita ut in quovis huius limitis puncto
unicum planum superficiem vasis tangens exstet. Si continuitas curvaturae in ali-
quo puncto singulari lineae P interrumpitur, sive cuspis ibi adsit, sive acies li-
neam P traiiciens, facile perspicietur, conclusionem nostram hinc non immutari;
sed aliter res se habet,    si continuitas curvaturae       interrupta   est in parte finita li-
neae P, i. e. si superficies vasis per partem finitam lineae P (vel adeo per totam
hanc lineam) aciem offert.        Tunc scilicet in quovis talis partis puncto bina plana
70                               PRINCIPIA   GENERALIA   THEORIAE


superficiem vasis tangentia aderunt, quorum alterum refertur ad partem liberam
superficiei vasis, alterum ad partem T. Betinendo itaque characterem   pro an-
gulo inter planum prius atque planum tangens superficiem U, denotandoque per
k angulum inter hoc planum et planum posterius, haud amplius erit i+k ---= 180\textdegree\{\},
sed' maior minorve, prout acies est convexa vel concaya. Et dum elementum va-
riationis    CW, pro motu virtuali fluidi ultra limitem P redundantis, etiamnum
exprimitur per

aadP.öe.sin(5, 6). (cos A---.cosi)

elementum illius variationis pro motu virtuali fluidi citra limitem P recedentis
iam erit

\[                           ---aadP.6e.sin(5, 6). (cos A+ cos k)\]

Ne   igitur valor ipsius    W capax sit variationis negativae, requiritur, ut neque va-
lor ipsius cos A---cosi        sit negativus, neque valor ipsius cos A+cosk             positivus,
i. e. esse debet

vel    i     A,               vel   :>A
atque vel    k---= 1800---A,            vel   k>180\textdegree\{\}----A

In statu aequilibrii itaque esse nequit ©+-\%k<\{180\textdegree\{\}, sive, quod idem est, in
statu aeqwlibrii limes superficiei fluidi liberae U esse nequit, per extensionem fini-
tam, in acie concava superficiei vasis.        Contra, quoties pars illius limitis coincidit
cum acie convexa, ad aequilibrium requiritur et sufficit, ut angulus inter plana
fiuidum et vas tangentia sit inter limites A et A-+-a(inel.), extra fluidum, sive
inter   180'---A et 180'--- A-+a,                intra fluidum mensuratus, si angulum inter
duo plana superficiem vasis utrimque ab acie tangentia in quovis puncto indefinite
per 180\textdegree\{\}'---a denotamus, quatenus hic angulus a plaga vasis mensuratur.


\subsection*{31.}

Constantes     aa, 55,    quarum ratio angulum          A      determinat, a functioni-
bus f, F pendent, et quodammodo tamquam mensurae intensitatis virium mo-
lecularium,    quas particulae fluidi et vasis exercent,            considerari possunt..     Si
functiones    istae ita comparatae      sunt, ut fx, Fx         sint    in ratione determinata
a distantia    x independente, puta ut           n ad N,       manifesto statuere possumus
FIGURAE   FLUIDORUM    IN STATU AEQUILIBRI.                    71


'aa:655 = cn:CN,         i.e, cohstantes   aa, 66 erunt proportionales attractionibus,
quas in eadem distantia exercent duae moleculae quoad volumen aequales, al-
tera fluidi altera vasis. Iam quum angulus A fiat acutus, recto aequalis, ob-
tusus, duobus rectis aequalis, prout bb<4aa, 55 --- taa, bb>+aa sed
\[<.aa, 656 = aa: in sensu istius suppositionis (quae si gratuita est, tamen veri-\]
similitudini non repugnat) dicere oportet, casum primum locum habere, quoties
attractio partium fluidi mutua maior sit quam duplum ättractionis partium vasis
in fluidum; secundum, quoties prior attractio sit duplum posterioris; tertium,
quoties prior maior quidem sit posteriori, sed minor eius duplo; denique quartum,
quoties ambae attractiones sint aequales.   Exemplum casus primi exhibet argen-
tum vivum in vasibus vitreis.



\subsection*{32.}

At quantus est valor anguli A in casu eo, ubi attractio vasis maior est
quam attractio partium fluidi mutua?   Valor imaginarius, quem pro 65 >aa
formula sin+A ---       angulo A assignat, iam testatur, suppositionem aliquam
in tali casu non admissibilem subesse.           Revera quoties    665 >aa,   suppositio Ä-
mitationis superficiei   7' cum conditione minimi respectu functionis         W consistere
nequit.    Ubicunque enim limitem posueris, patet, si ultra hunc limitem cutem
fluidi tenuissimam expansam concipias, ita ut 7' capiat augmentum              T, et proin
U augmentum huic proxime aequale, valorem functionis W assumere mutatio-
nem sensibiliter aequalem quantitati negativae --- (26565---2aa)T'; quinadeo va-
lorem ipsius W tamdiu ulterioris diminutionis capacem manere, donec 7' totam
superficiem vasis religquam occupaverit. Valor mutationis --- (255 ---2aa)T' eo
magis exactus erit, quo minor crassities accipiatur, et quatenus tantummodo de
valore expressionis W agitur, nihil impedit, quominus crassities usque ad eva-
nescendum diminui concipiatur.         Attamen cutis crassitiei evanescentis (probe di-
stinguendae ab insensibili) nihil esset nisi fictio mathematica,        figuraque spatü s
tali fictioni accommodata revera non differret ab ea, pro qua W in casu bb =aa
valorem minimum acquirit.
Sed paullo aliter res se habet in problemate nostro physico, ubi talis cutis
accessoria necessario gaudere debet certa crassitie, utut insensibili, quo aequili-
brium consistere possit. Quoties talis pars adest, expressio W, uti in art. 18 do-
cuimus, incompleta est, et denotata ea parte vasis, quam cutis tegit, per 7'',
72                            PRINCIPIA GENERALIA THEORIAE

huiusque crassitie in quovis puncto indefinite per p,* expressioni @ adhuc adii-
ciendi erunt termini

\[                              rec/hp.dT'---\]
rel [Op.AT'

adeoque valori ipsius W hi
rc       '          r       rc               r



r       ,256       ,        200    nr




Quocirca quum valor ipsius      W, per accessionem istius cutis, iam acceperit mu-
tationem   (2656---2aa)T',       mutatio tota, ei valori ipsius                        W,   qui omittendo cu-
tem locum habet,     adiicienda , erit

\[                                   1--- gt)---aa(\i\{\}\]
\[                           --- 2[aT.[66                                            42]\]
Haec mutatio propter      0'0 = \#0, 8'0 = 00,                        nulla esset pro crassitie evane-
scente; at quum D'p, 8'p, crescente crassitie p, citissime decrescant, et iam pro
valore insensibili ipsius p insensibiles evadant, mutatio ista citissime versus va-
lorem ---       (2656
--- 2aa)T" converget, atque pro statu aequilibrii fluidi, ne va-
lor expressionis   W correctae capax sit ulterioris diminutionis sensibilis, sensibi-
liter eidem aequalis esse debebit.    Ceterum investigatio completa legis, quam
crassities p sequi debet, profundiores evolutiones requireret, quibus tamen hic
non immoramur, quum absque cognitione functionum f, F, a quibus functiones
6', 0\textdegree\{\} pendent, nec non propter rationes in art. 34 indicandas, nimis otiosae vi-
deri possent. Ad investigationem partis substantialis fluidi, i. e. eius, cuius di-
mensiones omnes sensibiles sunt, sufficit, pro casu nostro, ubi 665>aao, vasin
vicinia limitis partis substantialis madefactum concipere, i.e. cute fluida obductam,
cuius crassities insensibilis quidem sit, attamen tanta, ut 0'p, 0'\% negligi pos-
sint. Hoc pacto functio, quae in statu aequilibrii minimum esse debet, erit
\[                     Szds--- 2(665---ao)(T+T)---aaT+aaU\]
ubi 7, U ad solam partem substantialem fluidi referri supponuntur.                                Patet ita-
que, variationem huius functionis e mutatione virtuali figurae partis substantialis
. Auidi oriundam (qualis mutatio aggregatum                       T-+-T' non afficit) convenire cum
variatione expressionis


fzas--- aa TtaaU

i. e. eiusdem expressionis, quae minimum esse debet pro casu 66--- aa.       Hinc
colligimus, figuram aequilibrii fluidi in vase, pro quo 56 >aca, convenire cum
figura aequilibrii eiusdem fluidi in vase, pro quo 665 = \guillemotleft\{\}aa, ea tamen differen-
tia, ut illa in aequilibrio stricto desinere debeat in cutem crassitiei insensibilis.
Ceterum ill. Lartace iam monuit, pro illo casu vas cute fluidi insensibilis cras-
sitiei obductum aequipollere vasi tali, cuius particulae vim attractivam in fluidum
exerceant vi attractivae partium fluidi mutuae aequalem.
Sponte hinc sequitur modificatio, propositionibus art. 18 circa ascensum flui-
dorum in tubis capillaribus verticalibus adiicienda: quoties scilicet 66>aa, in
formulis illie allatis \guillemotleft\{\}a loco ipsius 66 substituere oportet.


\subsection*{33.}

In casu eo, ubi     66 <\{a\guillemotleft\{\}a,       madefactio vasis per cutem fluidi insensibilis
crassitiei loecum habere nequit, siquidem lex functionum            6', 8' ea est, ut valor
functionis

\[                               aal\i\{\} \%,_ BB             --- en)\]

pro qua brevitatis   caussa   scribemus     Qp,    continuo   crescat,   dum    p a valore    0
versus valorem sensibilem progreditur: manifesto enim pro tali functionis             Qp in-
dole existentia talis cutis conditioni minimi repugnaret.           Sponte illam indolem
affert hypothesis, de qua in art. 31 loquuti sumus, puta ubi /x, Fx sunt in ra-
tione determinata ab x independente, quoniam hinc etiam sequitur;- = ne et
proin Qp = (na---56)(1 .                At si functiones f, F legem van                 seque-
rentur, haud impossibile esset, ut valore ipsius ni rapidius decrescente, quam
valore ipsius a functio Qp, intra ambitum valorum insensibilium ipsius p,
primo fieret negativa, et postquam attigisset valorem suum minimum (i. e. extre-
mum negativum) rursus ascenderet per valorem 0 versus limitem suum positivum
aa--- 66. In tali casu aequilibrium utique postularet eutem insensibilem, euius
crassities generaliter loquendo tanta esse deberet, ut Qp haud sensibiliter discre-
peta valore suo minimo. Qui si per ---6'6' denotatur, erit 655 <66; figura
autem partis substantialis fuidi perinde determinabitur, ac si esset in vase, cuius
respectu loco quantitatis 66 substituere oportet 6'5', i.e. angulus inter planum
2                                                                        10
74                                     PRINCIPIA   GENERALIA          THEORTAE


superficiem fluidi liberam in confiniis partis substantialis tangens atque parietem
vasis erit ---         2arc.sin,.         Sed quum         valde dubium sit, an talis casus in rerum
*        .             .    6'                                      .        [2          *         [3




natura exstet,         superfluum videtur,          diutius ei immorari.



\subsection*{34.}

Alienum foret ab instituto nostro praesente, a principiis generalibus hic sta-
bilitis ad phaenomena specialia descendere, praesertim quod illorum principiorum
essentia quadrat cum theoria ea, per quam ill. Larrace aequali arte et successu
permulta phaenomena in aequilibrio fluidorum conspicua iam explicavit. Vastus
utique superestcampus, largam messem novam pollicens: sed haec curis futuris
reservata maneat.  Contra e re erit, quasdam annotationes adiicere, quae vel
novam lucem huic argumento affundere, vel interpretationem erroneam arcere
poterunt.


I.       Theoria nostra non arrogat sibi'determinationem figurae aequilibrii ma-
thematice exactam,              sed acquiescit in determinatione                  figurae talis, a qua figura
aequilibrii vera differre nequit quantitate sensibili. Errares, si hoc alicui imper-
fectioni theoriae tribueres, quae ex asse praestitit, quantum praestare possibile
est, quamdiu lex attractionis molecularis ignoratur. In statu aequilibrii functio
Q exacte maximum esse debet, adeoque functio
2rcsbo        Q
g           ge

minimum;           haec autem,        pro indole attractionis molecularis,                  non quidem exacte
aequalis est functioni            W,     attamen insensibiliter tantum ab ea differt.                        Figura
igitur, pro qua W fit minimum, non est exacta figura aequilibrii, sed differen-
tia esse debet insensibilis, quatenus quidem quaelibet mutatio sensibilis istius
figurae valorem sensibiliter maiorem functionis                          W produceret.            Manifesto' hinc
non excluditur differentia sensibilis in curvatura                           superficiei,   dummodo     limitetur
ad partem superficiei insensibilem: quapropter in figura aequilibrii exacta angu-
lum constantem supra per A denotatum haud amplius considerare licet tamquam
inclinationem superficiei fluidi ad parietem vasis in ipso contactu, sed tantum-
modo in distantia immensurabili                a vase,         sive, ut ill. LarzaceE recte iam monuit,
inclinatio in limite sphaerae sensibilis attractionis vasis cum valore ipsius A sen-
sibiliter coincidet.
FIGURAE FLUIDORUM IN STATU AEQUILIBRI.                     75

II. Probe distinguere oportet figuram aequilibrii a figura quietis.   Quoties
fluidum est in statu aequilibrii, certo in eo perseverare debebit. At quoties figura
fluidi aliquantum a figura aequilibrii differt, nihilominus accidere potest, ut flui-
dum vel in quiete permaneat, vel, si moveatur, motum iam amittat, antequam
statum aequilibrüi attigerit, perinde ut e. g. cubus plano horizontali tantum im-
positus in aequilibrio versatur, sed etiam supra planum inclinatum quiescere pot-
est, frictione motum impediente.     Ita fluidum talem statum occupans, pro quo
W habet valorem minimum, certo quiescet: sed quoties est in statu ab illo di-
verso, puta ubi       W diminutionis capax est, ex hoc statu in statum aequilibrii ea-
tenus tantum transibit, quatenus frictio non impediverit.           Hocce autem respectu
duae conditiones aequilibrii essentialiter diversae sunt.       Scilicet aequatio funda-
mentalis prior (art. 28) independens est a mutabilitate limitis P, i.e. ad condi-
tionem minimi tunc quoque necessaria, ubi hic limes invariabilis supponitur: qua-
propter, quatenus quidem fluidum perfecta fluiditate gaudet, ut pars una supra
alteram libere gliscere possit, dum vel minima vis motum postulat, fluidum ne-
cessario illi conditioni se accommodabit.   Longe vero alia est ratio principii se-
cundi (art. 29), quod essentialiter pendet a perfecta limitis P mobilitate in super-
ficie vasis. Conditio minimi in valore ipsius W utique postulat aequationem
i= A: si vero, postquam superficies fluidi priori quidem prineipio se accommo-
davit, angulus © nondum assequutus est valorem normalem, neque adeo W va-
lorem absolute minimum, transitus in statum aequilibrii perfecti fieri nequit abs-
que translocatione limitis P, sive absque motu fluidiin contactu cum vase, quali
motui utique obstare potest frictio.  Hinc manifestum est, cur in experimentis
circa eadem corpora institutis tantas differentias in valore anguli ö offendamus.
Perinde in casu eo, ubi       656>oaa,    fluidum in vase, cuius parietes iam sunt ma-
defacti, utique se componet ad legem aequilibrii, secundum quam pro parte sub-
stantiali fluidi esse debet = 180\textdegree\{\}: sed in vase, cuius parietes extra fluidum
etiamnum sunt sicci, fluidum a statu non aequilibrato proficiscens parietesque va-
sis siccas invadens iam ad quietem pervenire poterit, antequam angulus i valo-
rem   180\textdegree\{\} attigit.     Hinc simul elucet ratio, cur phaenomena capillaria fluidorum
talium, quae madefactioni non adversantur, in tubis siccis tantas irregularitates
offerant, ascensumque saepissime longe minorem, quam in tubis iam humectatis,
ubi consensum pulcherrimum cum theoria semper aspicimus.

10*
76                                     PRINCIPIA   GENERALIA       THEORIAE


III. : Ratio constantium               \guillemotleft\{\}\&,6 e phaenomenis derivari nequit, quoties 6
est maior quam a: figura enim eiusdem fluidi in vasibus forma aequalibus ma-
teria diversis pro isto casu non differt nisi respectu cutis immensurabilis vas ma-
defacientis.   Quoties autem Ö minor estquam \guillemotleft\{\}a, determinatio rationis inter has
constantes possibilis quidem est adiumento anguli i, sed propter rationes modo
allatas magnam praecisionem vix feret. Pro mercurio in vasibus vitreis ill. La-
PLACE statuit angulum i = 43\textdegree\{\} 1?..
Longe maioris praecisionis capax est determinatio constantis                           a, praesertim
si vasibus madefactionem admittentibus uti licet. Pro aqua sub temperatura
8,5 graduum thermometri centesimalis statuere oportet secundum experimenta
ab ill. Lapuace citata*)
aa   --- 7,5675 millim. quadr., sive
a ---     2,7509 millim.

Pro spiritu vini, cuius pondus specificum = 0,81961,                           sub eadem temperatura
oa --- 3,0441 millim. quadr., sive
a. |--- 1,7447 millim.                        |
Pro oleo terebinthino sub temperatura                    8 graduum

ac ---= 3,305 millim. quadr., sive
a\guillemotleft\{\} ---= 1,818 millim.

Pro mercurio, sub temperatura 10 graduum,                            statuere licet, donec experimenta
nova maiorem praecisionem suppeditaverint,

aa ---= 3,25 millim. quadr., sive
\[                                    a = 1,803 millim.\]

Ceterum verisimile est, temperaturam eatenus tantum valorem constantis aa af-
ficere, quatenus densitas inde pendet, cui itaque in hac hypothesi valor ipsius
aa proportionalis erit.
Er


\footnote{Observare   convenit,    quantitatem   ab ill. Larzace   per H denotatam   convenire cum nostro   r.c\#o,}

adeoque   =   apud illum auctorem     idem denotare,   quod in signis nostris est an ke      Te
FIGURAE   FLUIDORUM   IN STATU   AEQUILIBRIL,                   77


Valores isti conclusi sunt ex ascensione vel depressione fluidorum in tubis
. capillaribus: attamen valde difficile est, horum diametros exacte mensurare, dif-
\_fieilius, de forma circulari sectionis transversalis certitudinem acquirere. Longe
maiorem praecisionem pollicentur experimenta circa diametros et volumina magna-
rum guttarum mercurii tabulae horizontali vel curvaturae perparvae notae insi-
dentium, qualia iam instituerunt physici SEGNErR et Gay-Lussac: nec non, pro
liquidis vasa   vitrea   madefacientibus,      experimenta     circa dimensiones   bullarum
magnarum aeris in vasibus superne operculo madefacto plano horizontali vel pa-
rum et secundum radium notum curvato clausis, ad quae instituenda physicos
invitamus.


|   IV. Ne limites huic commentationi praescriptos excederemus, applicatio-
nem principiorum nostrorum generalium hocce quidem loco ad casum simplicis-
simum restringere oportuit, ubi liquidum unicum in vase firmo consideratur.
Nihil vero obstat, quominus theoriae summa generalitas concilietur, ita ut etiam
problema plurium liquidorum in eodem vase, nec non casum eum amplectatur,
ubi insuper corpora rigida, vel omnino vel ex parte libera, fluido immersa sunt.
Sed harum quaestionum uberiorem expositionem ad aliam occasionem nobis re-
servare debemus.
INTENSITAS
VIS MAGNETICAE                                   TERRESTRIS
AD MENSURAM                 ABSOLUTAM          REVOCATA.











IN CONSESSU    SOCIETATIS   MDCCCXXXII.    DEC. XV. RECITATA.




Commentationes societatis regiae seientiarum Gottingensis recentiores.   Vol. v\i\{\}n.
INTENSITAS
VISMAGNETICAE TERRESTRIS
AD MENSURAM    ABSOLUTAM     REVOCATA.



Ad determinationem completam vis magneticae telluris in loco dato
tria ele-
menta requiruntur: declinatio seu angulus inter planum in quo agit atque
planum
meridianum; inclinatio directionis ad planum horizontale; deniqu
e tertio loco in-
tensitas. Declinatio, quae respectu omnium applicationum ad
usus nauticos at-
que geodaeticos tamquam elementum primarium consideranda
est, statim ab ini-
tio observatores atque physicos exercuit, qui etiam inclinationi
curas assiduas iam
per saeculum         dicaverunt.   Contra   elementum
tertium, intensitas, licet aeque
dignum scientiae obiectum, usque ad tempora recentiora
penitus neglectum man-
sit. Illustri Humsorpr inter tot alias ea quoque laus debetu
r, quod primus fere
huic argumento animum advertit, inque itineribus suis
magnam copiam observa-
tionum circa intensitatem relativam magnetismi terrestris congess
it, e quibus con-
tinuum incrementum huius intensitatis, dum ab aequatore
magnetico versus po-
lum progredimur, innotuit. Permulti physici, vestigiis huius naturae scrutatoris
insistentes,       iam tantam determinationum    copiam    contulerunt,   ut clarissimus et
de magnetismo terrestri meritissimus vir, HANsTerN, specim
en mappae universa-
, lis lineas isodynamicas exhibentis nuper iam edere potuerit.
Methodus, qua in hoc negotio utuntur, consistit in observ
atione temporis,
,per quod eadem acus magnetica in locis diversis numerum
oscillationum eundem
perficit, seu numeri oscillationum in temporis intervallo eodem
, intensitasque
quadrato numeri oscillationum in tempore dato proportionalis
ponitur: hoc modo
inter se comparantur intensitates totales, dum acus inclinatoria
in centro gravi-
tatis suspensa circa axem horizontalem ad meridianum magne
ticum normalem
oscillat, seu intensitates vis horizontalis,    dum     acus horizontalis circa axem ver-
V.           .
11
82                       INTENSITAS   VIS MAGNETICAE   TERRESTRIS


ticalem vibratur: posterior observandi modus maiorem praecisionem fert, et quae
inde resultant, cognitis inclinationibus, facile ad intensitates totales reducuntur.
Manifesto nervus huius methodi pendet a suppositione, distributionem
magnetismi liberi in particulis acus ad talem comparationem adhibitae in singulis
experimentis invariatam mansisse: si enim vis magnetica acus lapsu temporis ali-
quantulam debilitationem passa esset, ob id ipsum postea tardius oscillaret, ob-
servatorque talis mutationis inscius intensitati magnetismi terrestris pro loco po-
steriori valorem nimis parvum tribueret.    Quodsi experimenta temporis interval-
lum mediocre complectuntur, acusque ex chalybe bene durato confecta et magne-
tismo caute imbuta in usum vocatur, considerabilis vigoris debilitatio parum uti-
que metuenda erit; praeterea incertitudo minuetur, si plures acus ad comparatio-
nes adhibentur; denique isti suppositioni maior fides accrescet, si peracto itinere
acus in loco' primo tempus vibrationis non mutasse invenitur.      Sed quaecunque
cautelae adhibeantur, lenta aliqua debilitatio vis acus vix evitari, adeoque talis
consensus post longiorem absentiam raro exspectari poterit; quapropter in com-
paratione intensitatum pro locis terrae valde dissitis plerumque tantam praecisio-
nem et certitudinem,   quantam desiderare debemus,         attingere haud licebit.
Ceterum hoc methodi incommodum minus grave est, quamdiu tantum de
comparatione intensitatum simultanearum vel temporibus non longe inter se di-
stantibus respondentium agitur. At quum experientia docuerit, tum declinatio-
nem   tum inclinationem in loco dato mutationes continuas pati, quae post multos
annos pergrandes evadant, dubium esse nequit, quin intensitas quoque magne-
tismi terrestris similibus mutationibus quasi saecularibus obnoxia sit; manifesto
autem, quatenus de hac quaestione agitur, methodus ista prorsus inefficax eva-
dit.    Et tamen, ad scientiae naturalis incrementum summopere desiderandum
foret, ut haec ipsa quaestio gravissima in plenissimam lucem promoveatur, quod
certo fieri nequit, nisi methodo pure comparativa abrogata alia substituitur, quae
a fortuitis acuum inaequalitatibus prorsus independens intensitatem magnetismi
terrestris ad unitates stabiles mensurasque absolutas revocat.
Haud difficile est, principia theoretica stabilire, quibus talis methodus, diu
iam in votis habita, innitidebet.     Multitudo oscillationum,
quas acus in tempore
dato perficit, pendet tum ab intensitate magnetismi terrestris, tum a constitutione
acus, puta a momento statico elementorum magnetismi liberi in illa contentorum,
atque ab eius momento inertiae.e    Quum hoc momentum inertiae absque difficul-
AD   MENSURAM   ABSOLUTAM   REVOCATA.                   83


tate assignari possit, patet, observationem oscillationum nobis suppeditare pro-
ductum ex intensitate magnetismi terrestris in momentum staticum magnetismi
acus: sed hae duae quantitates separari nequeunt, nisi observationibus alius ge-
neris accitis, quae diversam earum combinationem implicant.       Ad hunc finem
accedat acus secunda, quae exponatur actioni et magnetismi terrestris et magne-
tismi acus primae, ut, quam rationem inter se teneant hae duae actiones, explo-
rari possit. Utraque quidem actio pendebit a distributione magnetismi liberi in
secunda acu: sed posterior insuper a constitutione acus primae, distantia centro-
rum, positione rectae centra iungentis respectu axium magneticorum utriusque
acus, denique a lege, quam attractiones et repulsiones magneticae sequuntur.
Immortalis Top\i\{\}as MAYER primus iam coniectaverat, hanc legem cum lege gravi-
tationis eatenus convenire, quod illae quoque actiones decrescant in ratione du-
plicata distantiarum: experimenta clarissimorum virorum CouLoms et HAnsTEEN
magnam huic coniecturae plausibilitatem conciliaverunt, experimentaque novis-
sima eam ultra omne dubium evehunt.          Sed probe attendendum est, hanc legem
referri ad singula elementa magnetismi liberi: effectus totalis corporis magnetismo
imbuti longe aliter se habebit, atque in distantiis permagnis, uti ex illa ipsa lege
deducere licet, proxime ad rationem inversam triplicatam distantiarum accedet, ita
ut actio acus per cubum distantiae multiplicata, distantia ceteris paribus continuo
erescente, ad valorem constantem asymptotice convergat, qui, dum distantiae,
linea arbitraria pro unitate accepta, per numeros exprimuntur, cum actione vis
terrestris homogeneus atque comparabilis erit. Per idoneam experimentorum ad-
ornationem et tractationem limes huius rationis eruendus est; qui quum     tantum-
modo momentum staticum magnetismi primae acus involvat, iam habebitur quo-
tiens e divisione huius momenti per intensitatem vis terrestris ortus, qui compara-
tus cum producto harum quantitatum antea eruto eliminationi istius momenti sta-
tici inserviet, atque valorem intensitatis magnetismi terrestris suppeditabit.
Quod attinet ad modum, actiones magnetismi terrestris et acus primae in
acum secundam ad experimenta revocandi, duplex via patet, quum acum secun-
dam vel in statu motus vel in statu aequilibrii observare possimus.     Prior modus
eo redit, ut oscillationes huius acus observentur, dum actio magnetismi terrestris
coniungitur cum actione acus primae in distantia idonea ita collocatae, ut ipsius
axis sit in meridiano magnetico per centrum acus oscillantis ducto: hoc pacto os-
eillationes vel accelerabuntur vel retardabuntur, prout poli amici vel inimici sibi
2 Aa

\section*{INTENSITAS VIS MAGNETICAE TERRESTRIS AD MENSURAM ABSOLUTAM REVOCATA.}
\addcontentsline{toc}{section}{Intensitas vis magneticae terrestris ad mensuram absolutam revocata.}


mutuo obversi sunt, comparatioque vel temporum vibrationum pro utraque acus
primae positione inter se, vel temporis alterutrius cum tempore vibrationum sub
sola magnetismi terrestris actione (remota acu prima) locum habentium, rationem
huius vis ad actionem primae acus docebit. Alter modus acum primam ita col-
locat, ut directio vis quam in regione acus secundae libere suspensae exercet, fa-
ciat angulum (e. g. rectum) cum meridiano magnetico, quo pacto haec ipsa a me-
ridiano magnetico deflectetur, et e magnitudine deflexus ratio inter vim magne-
ticam terrestrem atque actionem acus primae concludetur.
Ceterum modus prior essentialiter convenit cum eo, quem ill. Po\i\{\}ssox iam
ante aliguot annos proposuit. Sed experimenta ad ipsius normam a nonnullis phy-
sicis tentata, quae quidem mihi innotuerunt,     vel successu omnino caruerunt,     vel
rudem tantummodo approximationem praebuerunt.
Difficultas rei inde potissimum pendet, quod ex actionibus acus in distantiis
mediocribus observatis computari debet limes aliquis, qui ad distantiam quasi in-
finite magnam refertur, et quod eliminationes ad hunc finem necessariae tanto ma-
gis a levissimis observationum erroribus turbantur, imo pervertuntur, quo plures
incognitae a statu acuum individuali pendentes eliminandae sunt: ad multitudi-
nem perparvam incognitarum autem res tunc tantummodo deduci potest, ubi actio-
nes in distantiis satis magnis (respectu longitudinis acuum) observatae sunt, adeo-
que ipsae iam perparvae evaserunt.      Sed ad actiones tam parvas accurate mensu-
randas subsidia practica hactenus usitata non sufliciunt.
Ante omnia itaque in id mihi incumbendum        esse vidi, ut subsidia nova pa-
rarem, per quae tum tempora oscillationum, tum directiones acuum longe maiori
praecisione, quam hactenus licuit, observare ac metiri possem.  Labores ad hunc
finem suscepti et per plures menses continuati, in quibus a praestantissimo WE-
BER multifariam adiutus sum, ad scopum exoptatum ita perduxerunt, ut exspecta-
tionem non modo non fefellerint, sed longe superaverint, nec iam quidquam de-
siderandum   restet,   ad praecisionem   experimentorum    subtilitati   observationum
astronomicarum 'aequiparandam,     nisi locus ab influxu ferri propinqui atque agi-
tationum aöris plene securus. Adsunt duo apparatus simplicitate non minus quam
praecisione quam praebent insignes, quorum discriptionem quidem ad aliam occa-
sionem mihi reservare debeo, dum experimenta ad determinandam intensitatem
magnetismi terrestris hactenus in observatorio nostro instituta physicis in hac com-
mentatione trado.
L
Ad explicationem phaenomenorum magneticorum duo fluida magnetica po-
stulamus; alterum cum physicis vocamus boreale, alterum australe. Elementa
fluidi alterius attrahere alterius elementa,    contra bina elementa eiusdem fluidi
mutuo se repellere supponimus, et quidem utramque actionem variari in ratione
inversa quadrati distantiae. Veritatem huius legis per ipsas nostras observatio-
nes plenissime confirmari infra apparebit.
Fluida ista non. per se apparent, sed tantummodo iuncta cum particulis pon-
derabilibus corporum talium, quae magnetismi capaces sunt, illorumque actiones
in eo se manifestant,   quod has vel ad motum    sollicitant,   vel motum,   quem aliae
vires in ipsas agentes, e. g. gravitas, producerent, impediunt vel mutant.
Actio itaque quantitatis datae fluidi magnetici in quantitatem datam vel
eiusdem fluidi vel alterius in distantia data comparabilis erit cum vi motrice data.
i. e. cum actione vis acceleratricis datae in massam datam, et quum fluida magne-
tica ipsa non nisi per effectus quos producunt cognoscere liceat, hi ipsi illorum
mensurae inservire debent.               |
Quo igitur hanc mensuram ad notiones distinctas revocare possimus, ante
omnia circa tria quantitatum genera unitates stabilire oportet,. puta unitatem
distantiarum, unitatem massarum ponderabilium, unitatem virium acceleratricium.
Pro tertia accipi potest gravitas in loco observationis : quod si minus arridet, in-
super accedere debet unitas temporis, eritque nobis vis acceleratrix ea =1,
| quae in unitate temporis mutationem velocitatis corporis in ipsius directione moti
unitati aequalem gignit.
86                              INTENSITAS   VIS    MAGNETICAE   TERRESTRIS


His ita intellectis,     unitas quantitatis fluidi borealis ea erit, cuius vis repul-
siva in aliam ipsi aequalem in distantia = 1 positam aequivalet vi motriei --- 1,
i. e. actioni vis acceleratricis ---= I in massam = 1, idemque de unitate quan-
titatis fluidi australis valebit: in hac determinatione manifesto tum fluidum agens,
tum fluidum in quod agitur, in punctis physicis concentrata concipi debent.  In-
super autem supponere oportet, attractionem inter quantitates datas fluidorum
heteronymorum in distantia data aequalem esse repulsioni inter quantitates resp.
aequales fuidorum homonymorum.                      Actio itaque quantitatis m fluidi magnetici
borealis in quantitatem m eiusdem fluidi in distantia r (dum utrumque in puncto
s                                   i             a              mm!         .       ;          ern     2    mm
physico concentratum supponitur),                  exprimitur per     ---, , Sive vi motricl                   ---
in directione a priori versus           posterius     agenti aequivalet,         manifestoque              haec for-
mula generaliter valet, si, quod semper abhinc subintelligemus,                               quantitas fluidi
R                        A                        5                 R        R       mm'
australis tamquam       negativa spectatur, ubi valor negativus vis                           ---,     Pro repul-
sione attractionem indicabit.                                     a
S\i\{\} itaque in puncto physico aequales quantitates fluidi borealis et australis
simul adsunt,      nulla omnino         actio hinc orietur,      si vero inaequales,                excessus al-
terius tantum, quem magnetismum \&berum (positivum seu negativum) vocabimus,
in considerationem veniet.



\subsection*{2.}

Hisce suppositionibus fundamentalibus adhuc aliam, quam experientia un-
dique confirmat, adiicere oportet, scilicet, quodvis corpus, in quo fluida magne-
tica adsint, semper aequalem utriusque quantitatem continere. Quinadeo expe-
rientia docet, hancce suppositionem etiam ad singulas talis corporis partes quan-
tumvis parvas, dummodo                sensibus     nostris discerni possint, extendendam                         esse.
Sed quum per ea, quae in fine art. praec. monuimus,                         actio eatenus tantum exi-
stere possit, quatenus aliqua separatio fluidorum locum habet, necessario hanc
per intervalla tam parva fieri supponere debemus, ut mensuris nostris non sint
accessibilia.
Corpus itaque magnetismi capax concipi debet tamquam compages innume-
rarum particularum,       quarum quaevis certam quantitatem                      fluidi magnetici borea-
lis et aequalem australis contineat, ita quidem, ut vel uniformiter inter se mixtae
sint (magnetismus lateat), vel separationem minorem maioremve iniverint (magne-
tismus evolutus sit), quae tamen separatio numquam                      in transfusionem fluidi ab
AD MENSURAM   ABSOLUTAM   REVOCATA.                         87


una particula in aliam abire potest. Nihil refert, utrum separatio maior a maiori
quantitate fluidorum quae libera evaserunt, an a maiori intervallo interposito orta
supponatur: manifesto autem praeter magnitudinem separationis simul eius di-
rectio in considerationem venire debet, quae prout in diversis corporis particulis
vel conspirat vel refragatur, maior minorve energia totalis respectu punctorum
extra corpus oriri poterit.
Quomodocunque autem distributio magnetismi liberi intra corpus se habeat,
semper eius loco, per theorema generalius, substituere licet secundum certam
legem aliam distributionem in sola corporis superficie, quae respectu virium ex-
trorsum agentium illi exacte aequivaleat, ita ut elementum fluidi magnetici extra
corpus ubicungue positum prorsus eandem attractionem vel repulsionem experia-
tur a distributione magnetismi vera intra corpus atque a fictitia in eius superficie.
Eandem fictionem ad bina corpora, quae ratione magnetismi liberi in ipsis evo-
luti in se invicem agunt, extendere licet, ita ut pro utroque distributio fictitia in
superficie distributionis verae internae vice fungi possit. Hocce demum modo vul-
gari loquendi mori, qui e.g. alteri acus magneticae extremitati solum magnetis-
mum borealem, alteri australem tribuit, sensum verum conciliare possumus, quum
manifesto haec phrasis cum principio fundamentali supra enunciato, quod alia
phaenomena imperiose postulant, non quadret.     Sed haec obiter hic annotavisse
sufficiat; de theoremate ipso, quum ad institutum praesens non sit necessarium,
alia occasione copiosius agemus.


\subsection*{3.}

Status magneticus corporis consistit in ratione distributionis magnetismi li-
beri in singulis eius particulis.   Respectu mutabilitatis huius status discrimen es-
sentiale inter corpora diversa magnetismi capacia animadvertimus.         In aliis, e.g.
in ferro molli, ille status per levissimam vim protinus mutatur, hacque cessante
status anterior redit: contra in aliis, praesertim in chalybe durato, vis certam
intensitatem attigisse debet, antequam sensibilem status magnetici mutationem
producere possit, vique cessante corpus vel in statu quem acquisivit permanet,
vel saltem ad priorem non ex asse'revenit.   In corporibus itaque prioribus mole-
culae fluidi magnetici semper ad aequilibrium perfectum virium, quae tum inter
ipsa mutuo,   tum a caussis externis     emanant,    se componunt,   vel saltem   a tali
aequilibrio sensibiliter vix differunt: contra in corporibus posterioris generis sta-
8383                      INTENSITAS VIS MAGNETICAE   TERRESTRIS

tus magneticus etiam absque perfecto aequilibrio inter illas vires durabilis esse
potest, si modo vires fortiores extraneae inde arceantur.            Etiamsi caussa huius
phaenomeni ignota sit, tamen eam ita imaginari licet, ac si partes ponderabiles
corporis secundi generis motui fluidorum magneticorum cum ipsis iunctorum ali-
quod obstaculum frictioni simile opponant, quae resistentia in ferro molli vel
nulla est, vel saltem perparva.
In disquisitione theoretica hi duo casus tractationem prorsus diversam re-
quirunt, sed in commentatione praesente de solis corporibus secundi generis sermo
erit: in experimentis, de quibus agemus, stabilitas status magnetici in singulis
corporibus ad illa adhibitis erit suppositio fundamentalis, probeque proin caven-
dum est, ne inter experimenta alia corpora, quae hunc statum mutare possent,
nimis prope accedant.
Attamen exstat quaedam caussa mutationis, cui etiam corpora secundi ge-
neris obnoxia sunt, puta calor. Nimirum experientia docet, statum magneticum
corporis variari cum eius temperatura, caloremque auctum intensitatem magne-
tismi debilitare,   ita tamen,   ut nisi corpus   ultra modum      calefactum fuerit, cum
'priori temperatura prior quoque status magneticus redeat. Haec dependentia per
experimenta idonea determinanda est, et si operationes ad idem experimentum
pertinentes sub temperaturis inaequalibus institutae sunt, ante omnia ad eandem
temperaturam revocandae erunt.


\subsection*{4.}

Independenter a viribus magneticis, quas corpora singularia satis sibi vi-
cina in se mutuo exercere videmus, alia vis in fluida magnetica agit, quam quum
ubique terrarum se manifestet, ipsi globo terrestri tribuimus, atque magnetis-
mum terrestrem vocamus.     Duplici modo haec vis se exserit: corpora secundi ge-
neris, in quibus magnetismus evolutus est, siin centro gravitatis sustinentur, ad
directionem determinatam sollicitantur; contra in corporibus primi generis fluida
magnetica per istam vim sponte separantur, quae separatio, si corpora figurae ido-
neae eliguntur atque in positione idonea collocantur, persensibilis reddi potest.
Utrumque phaenomenon explicatur, vim illam ita concipiendo, ut fluidum magne-
ticum boreale in quovis loco versus certam directionem propellat, australe vero
aequali intensitate versus oppositam. Directio prior semper intelligitur, dum de
directione magnetismi terrestris loquimur, quae proin per inclinationem ad pla-

num horizontale atque declinationem plani verticalis, in quo agit, a plano meri-
diano determinatur: illud planum meridianum magneticum vocatur.               Intensitas au-
tem magnetismi terrestris per vim motricem, quam in unitatem fluidi magnetici
liberi exserit,   aestimanda est.
Haec vis non modo in diversis terrae locis diversa est, sed etiam in eodem
loco variabilis, tum per saecula et annos,       tum per anni aestates dieique horas.
Respectu directionis haec variabilitas dudum quidem nota fuit: sed respectu in-
tensitatis hactenus tantummodo per horas diei animadverti potuit, quum subsi-
diis ad longiora temporis intervalla aptis caruissemus. Huic defectui in posterum
reductio intensitatis ad mensuram absolutam remedium afferet.


\subsection*{5.}

Ut actio magnetismi terrestris in corpora magnetica secundi generis (qualia
semper abhinc subintelligenda sunt) calculo subiiciatur, concipiatur tale corpus
in partes infinite parvas divisum, sitque dm            elementum 'magnetismi liberi in
particula, cuius coordinatae respectu trium planorum inter se normalium et re-
spectu corporis fixorum denotentur per       @, y, z:    elementa fluidi australis nega-
tive accipi supponimus.      Ita primo patet, integrale      7 dm   per totum corpus col-
lectum (imo per quamlibet corporis partem mensurabilem) esse --- 0. Statuamus
fedm = X, [ydm = Y, f[zdm --- Z, quae quantitates vocari poterunt mo-
menta magnetismi liberi respectu trium planorum fundamentalium, sive respectu
\_ axium in ipsa normalium.   Quum denotante a quantitatem constantem arbitra-
riam, fiat f (---a)dm = X,            patet, momentum respectu axis dati pendere tan-
tummodo ab eius directione, non autem ab eius initio. Si per initium coordina-
tarum axem quartum ducimus, qui cum primariis faciat angulos A, B, C, mo-
mentum elementi dm respectu huius axis erit --- (\guillemotleft\{\}cos A+ycosB-+zcosC) dm,
adeogue momentum magnetismi liberi in toto corpore
= XcosA+YcosB+ZcosC=V
Statuatur

\[ VX\&X+-YY+ZZ)=M,                atque X=Mcosa,               Y=Mcos6,          Z = Mecosy\]
ducaturque axis quintus, qui cum tribus primariis faciat angulos a, 6, y, et cum
axi quarto angulum     w; unde quum constet esse
v.                                                                    12\textdegree\{\}
90                           INTENSITAS    VIS MAGNETICAE   TERRESTRIS



cosw --- cos A cosa---+cosBcos6b----cosC cosy,                  fit    V---= Mecosw

Hunc axem quintum simpliciter vocamus                 corporis avem magneticum,     eiusque di-
rectionem ad valorem positivum radicalis              YXX+YY-+ZZ)              referri supponi-
mus.  Si axis quartus cum hoc axe magnetico coincidit, momentum V ft =M,
quod manifesto inter omnia est maximum: momentum respectu cuiuslibet alius
axis invenitur,      multiplicando   hoc    momentum        maximum      (quod quoties ambigui-
tas non metuenda est, simpliciter momentum magnetismi vocari potest) per cosi-
num anguli inter hunc axem atque axem magneticum. Momentum respectu cu-
jusvis axis in axem magneticum normalis fit = 0, negativum vero respectu cu-
iusvis axis, qui cum axe magnetico angulum obtusum facit.
Axis itaque magneticus non est recta determinata, quum per punctum quod-
libet duci possit, sed tantummodo directio determinata, sive adsunt infinite multi
axes magnetici inter se parallel.           E quibus si aliquem ad lubitum eligimus, lon-
gitudinemque determinatam ipsi tribuimus, eius termini vocantur poli, alter au-
stralis, a quo, alter borealis, versus quem directio axis procedit.


\subsection*{6.}

Si in singulas fluidorum magneticorum particulas vis agit intensitate et di-
rectione constans, vis totalis in corpus inde resultans facile e principiis statieis
derivatur, quum in corporibus, quae hic consideramus, particulae illae fluiditatem
quasi amiserint, et cum corpore ponderabili massam unam rigidam sistant.                  Agat
in guamvis moleculam magneticam dm                vis motrix    --- Pdm       secundum directio-
nem D (ubi pro moleculis fluidi australis signum negativum iam per se directio-
nem oppositam implicat); sint A, B duo corporis puncta in directione axis magne-
tici jacentia, eorumque distantia --- r, positive accepta, dum axis magneticus
tendit ab A versus B: ita facile intelligitur, si viribus istis duae novae adiun-
gantur,    utraque    ---= Sn, et quarum altera agat in A secundum                directionem   D,
alterain B secundum directionem oppositam, inter omnes has vires aequilibrium
fore. Quapropter vires priores aequivalebunt duabus viribus --- ap,, quarum
altera in B secundum directionem D, altera in A secundum directionem oppo-
sitam agit, manifestoque hae duae vires in unam conflari nequeunt.
Si praeter vim P alia similis P' secundum directionem .D' in corporis fluida
magnetica agit, eius loco iterum duae aliae vel in eadem puncta A, B, vel ge-


\end{document}
